%% ----------------------------------------------------------------
%% Introduction.tex
%% ---------------------------------------------------------------- 
\chapter{Introduction}\label{chapter:introduction}

% The recent development of low-altitude hypersonic vehicle technology for both civilian and military applications has prompted renewed scientific interest in the field. The hypersonic regime is conventionally defined as flows at speeds above Mach 5, where flight vehicles encounter strong shock-wave compressions that cause an abrupt increase in gas temperature. This gives rise to a \textit{high-enthalpy} environment, in which molecular dissociation in the atmosphere leads to chemical and vibrational nonequilibrium \citep{Candler2019}. The intricate nature of high-enthalpy flows induces complex interactions between the vehicle surface and the surrounding flow, characterised by the dissipation of a large amount of fluid kinetic energy into thermal energy within the viscous boundary layer. This, in turn, causes an acute rise in peak temperature over small wall-normal distances, leading to steep temperature gradients and elevated surface heat fluxes. 

% The recent development of low-altitude hypersonic vehicles for both civilian and military applications has renewed scientific interest in high-speed aerothermodynamics. The hypersonic regime is conventionally defined as flows above Mach~5, where strong shock-wave compressions dramatically raise the flow’s stagnation temperature, creating a high-enthalpy environment that drives molecular dissociation and vibrational excitation, leading to thermochemical nonequilibrium~\citep{Candler2019}. Within the boundary layer, viscous dissipation converts kinetic energy into thermal energy, producing steep temperature gradients and elevated surface heat fluxes. From an engineering perspective, accurately predicting these effects — along with associated pressure and shear stress distributions — is crucial for determining aerodynamic performance and for sizing the \gls{TPS}, which must protect the structure from the resulting extreme thermal environment. \fref{figure:trajectory} compares the trajectories of various existing and experimental vehicles, including the National Aerospace Plane (NASP) \citep{DSB1988,DSB1992}, a representative \gls{HGV}~\citep{Rizvi2015} and the Shuttle Orbiter, illustrating how different flight paths generate varying thermal loads that the \gls{TPS} must be designed to withstand. The boundary layer along these trajectories is affected by real gas effects that modify temperature and velocity profiles, as well as transport properties such as viscosity, thermal conductivity, and species diffusion. Capturing these effects accurately requires models that go beyond perfect-gas assumptions.

% The heat flux on the vehicle’s surface depends strongly on whether the boundary layer flow is laminar or turbulent, with heating loads increasing by a factor of 2–10 for turbulent flow~\citep{White2005} compared to laminar flow. This introduces further complexity not only for the design and sizing of the \gls{TPS} but also for predicting the aerodynamic forces and moments required to maintain stability and control throughout the entire flight envelope \citep{Dalle2011}. Determining the location or extent of transition to turbulence is a critical design consideration, but it remains highly challenging to predict with confidence. Predicting hypersonic boundary-layer transition is challenging due to the vast parameter space, its high sensitivity to flow conditions, uncertainties in the disturbance environment, and the inherent complexity of the nonlinear processes that drive the final breakdown to turbulence. 

The recent development of low-altitude hypersonic vehicles for both civilian and military applications has renewed scientific interest in high-speed aerothermodynamics. The hypersonic regime is conventionally defined as flows above Mach~5, where strong shock-wave compressions dramatically raise the flow's stagnation temperature, creating a high-enthalpy environment that drives molecular dissociation and vibrational excitation, leading to thermochemical non-equilibrium~\citep{Candler2019}. Within the boundary layer, viscous dissipation converts kinetic energy into thermal energy, producing steep temperature gradients and elevated surface heat fluxes. From an engineering perspective, accurately predicting these effects---along with associated pressure and shear stress distributions---is crucial for determining aerodynamic performance and for sizing the \gls{TPS}, which must protect the structure from the resulting extreme thermal environment. \fref{figure:trajectory} compares the trajectories of various existing and experimental vehicles, including the National Aerospace Plane (NASP)~\citep{DSB1988,DSB1992}, a representative \gls{HGV}~\citep{Rizvi2015}, and the Space Shuttle Orbiter, illustrating how different flight paths generate varying thermal loads that the \gls{TPS} must be designed to withstand. The boundary layer along these trajectories is affected by real-gas effects that modify temperature and velocity profiles, as well as transport properties such as viscosity, thermal conductivity, and species diffusion. Capturing these effects accurately requires models that go beyond perfect-gas assumptions.

% The heat flux on the vehicle’s surface depends strongly on whether the boundary-layer flow is laminar or turbulent, with heating loads increasing by a factor of 2–10 for turbulent flow~\citep{White2005,VanDriest1952} compared to laminar flow. This adds further complexity not only for the design and sizing of the \gls{TPS} but also for predicting the aerodynamic forces and moments required to maintain stability and control throughout the entire flight envelope~\citep{Dalle2011}. The location and extent of transition to turbulence are therefore critical design considerations, yet remain highly uncertain due to the vast parameter space, the sensitivity of transition to flow conditions, uncertainties in the disturbance environment, and the inherent complexity of the non-linear processes that drive the final breakdown to turbulence. This accurate prediction of transition onset and breakdown location essential for ensuring that the \gls{TPS} is neither under- nor over-designed for the mission conditions. Consequently, the use of overly conservative transition models can lead to unnecessary weight penalties in the \gls{TPS} design.

% The heat flux on the vehicle’s surface depends strongly on whether the boundary-layer flow is laminar or turbulent, with heating loads increasing by a factor of 2–10 for turbulent flow~\citep{White2005,VanDriest1952} compared to laminar flow. This adds further complexity not only for the design and sizing of the \gls{TPS} but also for predicting the aerodynamic forces and moments required to maintain stability and control throughout the entire flight envelope~\citep{Dalle2011}. The location and extent of transition to turbulence are therefore critical design considerations, yet remain highly uncertain due to the vast parameter space, the sensitivity of transition to flow conditions, uncertainties in the disturbance environment, and the inherent complexity of the non-linear processes that drive the final breakdown to turbulence. Accurate prediction of transition onset and breakdown location is essential to ensure that the \gls{TPS} is neither under- nor over-designed for mission conditions. Consequently, the use of overly conservative transition models can lead to unnecessary weight penalties in the \gls{TPS} design. 
The heat flux on the vehicle’s surface depends strongly on whether the boundary-layer flow is laminar or turbulent, with heating loads increasing by a factor of 2–10 for turbulent flow~\citep{White2005,VanDriest1952} compared to laminar flow. This adds further complexity not only for the design and sizing of the \gls{TPS} but also for predicting the aerodynamic forces and moments required to maintain stability and control throughout the entire flight envelope~\citep{Dalle2011}. The location and extent of transition to turbulence are therefore critical design considerations, yet remain highly uncertain due to the vast parameter space, the sensitivity of transition to flow conditions — including Reynolds number, Mach number, surface geometry, freestream turbulence and acoustic noise, pressure gradients, wall roughness and porosity, boundary-layer cooling or heating, and mass injection or suction~\citep{Cebeci2001} — uncertainties in the disturbance environment, and the inherent complexity of the non-linear processes that drive the final breakdown to turbulence. Accurate prediction of transition onset and breakdown location is essential to ensure that the \gls{TPS} is neither under- nor over-designed for mission conditions. Consequently, the use of overly conservative transition models can lead to unnecessary weight penalties in the \gls{TPS} design.







% The heat flux on the vehicle’s surface depends strongly on whether the boundary-layer flow is laminar or turbulent, with heating loads increasing by a factor of 2–10 for turbulent flow~\citep{White2005,VanDriest1952} compared to laminar flow. This adds further complexity not only for the design and sizing of the \gls{TPS} but also for predicting the aerodynamic forces and moments required to maintain stability and control throughout the entire flight envelope~\citep{Dalle2011}. The location and extent of transition to turbulence are therefore critical design considerations, yet remain highly uncertain due to the vast parameter space, sensitivity to flow conditions, uncertainties in the disturbance environment, and the inherent complexity of the non-linear processes that drive the final breakdown to turbulence. Accurate prediction of transition onset and breakdown location is essential to ensure that the \gls{TPS} is neither under- nor over-designed for mission conditions, as overly conservative models can lead to unnecessary weight penalties.




% External aerodynamic flows are inherently susceptible to environmental and surface disturbances that can destabilise the laminar boundary layer and promote transition. Turbulent flows enhance energy and momentum mixing compared to laminar flows, resulting in steeper gradients and surface heat fluxes that can be several times higher~\citep{VanDriest1952}. This makes accurate prediction of transition onset and breakdown location essential for ensuring that the \gls{TPS} is neither under- nor over-designed for the mission conditions. Consequently, the use of overly conservative transition models can lead to unnecessary weight penalties in the \gls{TPS} design.


\begin{figure}[hbt!]
    \centering
    \fontsize{9}{10}\selectfont\includesvg[width=0.96\textwidth]{resources/figures/introduction/schematictrajectory.svg}
    \caption{Schematic of a small-sized waverider going at hypersonic speed in the Earth’s atmosphere experiencing real gas phenomena adapted from \citep{Anderson2019} with the HGV glide path imposed from \citep{Rizvi2015}.}
    \label{figure:trajectory}
\end{figure}


% Downstream of the leading edge, the intense thermal and pressure gradients generated by the shock pose additional challenges for \gls{TPS} design and vehicle integrity. While translational and rotational modes equilibrate rapidly, vibrational and chemical relaxation occur over longer timescales, resulting in a persistently non-uniform thermochemical state along the glide path.  At low altitudes and high speeds, elevated atmospheric density produces high Reynolds numbers, making the boundary layer highly susceptible to transition. Even small perturbations — such as \gls{TPS} ablation, stagnation-region receptivity, or freestream disturbances — can trigger laminar–turbulent transition, which substantially increases thermal loads on the surface. High stagnation enthalpies and real gas effects further complicate this process by altering boundary-layer stability. Since turbulence can dramatically amplify heating and structural loads, accurately predicting transition onset is essential for reliable \gls{TPS} design and vehicle survivability.




% For a representative \gls{HGV}, downstream of the leading edge, the intense thermal and pressure gradients generated by the shock induce strong nonequilibrium effects that persist within the boundary layer. While translational and rotational modes equilibrate rapidly, vibrational and chemical relaxation occur over longer timescales, resulting in a persistently non-uniform thermochemical state along the glide path. At low altitudes and high speeds, elevated atmospheric density produces high Reynolds numbers, making the boundary layer more susceptible to transition. Even small perturbations — such as \gls{TPS} ablation, stagnation-region receptivity, or freestream disturbances — can promote laminar–turbulent transition, substantially increasing surface thermal loads. High stagnation enthalpies and real gas effects further complicate this by altering boundary-layer stability characteristics \citep{Anderson2019}. Since turbulence can dramatically amplify both heating and aerodynamic loads, reliably predicting transition onset is essential for robust \gls{TPS} performance and overall vehicle survivability.

% Transition is influenced by a variety of flow features, including Reynolds number, Mach number, surface geometry, inflow turbulence and acoustic noise, pressure gradients, wall roughness and porosity, boundary layer cooling or heating, and mass injection or suction through the surface \citep{Cebeci2001}. For a representative \gls{HGV}, downstream of the leading edge, the intense thermal and pressure gradients generated by the shock induce strong nonequilibrium effects that persist within the boundary layer. While translational and rotational modes equilibrate rapidly, vibrational and chemical relaxation occur over longer timescales, resulting in a non-uniform thermochemical state along the glide path. At low altitudes and high speeds, elevated atmospheric density produces high Reynolds numbers, making the boundary layer more susceptible to transition. Even small perturbations — such as \gls{TPS} ablation, stagnation-region receptivity, or freestream disturbances — can promote transition and substantially increase surface thermal loads \citep{miromiroThesis}. High stagnation enthalpies and real gas effects further complicate this by altering boundary-layer stability characteristics~\citep{Anderson2019}. These coupled effects highlight the importance of reliable transition prediction to maintain both thermal protection and aerodynamic control without excessive design margins.

For a representative \gls{HGV}, downstream of the leading edge, the intense thermal and pressure gradients generated by the shock induce strong nonequilibrium effects that persist within the boundary layer. While translational and rotational modes equilibrate rapidly, vibrational and chemical relaxation occur over longer timescales, resulting in a non-uniform thermochemical state along the glide path. At low altitudes and high speeds, elevated atmospheric density produces high Reynolds numbers, making the boundary layer more susceptible to transition. Even small perturbations — such as \gls{TPS} ablation, stagnation-region receptivity, or freestream disturbances — can promote transition and substantially increase surface thermal loads~\citep{miromiroThesis}. High stagnation enthalpies and real gas effects further complicate this by altering boundary-layer stability characteristics~\citep{Anderson2019}. These coupled effects highlight the importance of reliable transition prediction to maintain both thermal protection and aerodynamic control without excessive design margins.


Accurately modelling the physics under such extreme conditions presents significant challenges, both experimentally and computationally~\citep{DiRenzo2021}. Only a few experimental facilities can replicate the high-enthalpy environments relevant to hypersonic flows, and even these are limited by the substantial power demands of test rigs and the need for quiet flow conditions in transition studies. As a result, most research has focused on the early stages of transition, while a comprehensive understanding of the breakdown pathways by which instability modes evolve into turbulence remains incomplete for hypersonic boundary layers. A deeper understanding of boundary-layer behaviour at high Mach numbers and low altitudes — where elevated temperatures, turbulence, and real gas effects interact — is crucial for safe and efficient vehicle design. These aspects remain insufficiently explored, highlighting the need for studies that address their coupled effects.






\section{High enthalpy flows}
% Vehicles operating in hypersonic flows encounter intense shock-wave compression, forming a high-enthalpy gas environment characterised by complex physicochemical processes. Under these extreme conditions, the coupling between macroscopic flow properties and the microscopic internal energy modes of the gas becomes significantly more pronounced, necessitating more advanced modelling than that used for classical, low-enthalpy flows \citep{Candler2019}. These internal energy modes — including translation, rotation, vibration, dissociation and exchange, electronic excitation, radiation, and ionisation — absorb, store, and release energy at discrete, quantised levels \cite[ch.~11]{Anderson2019}. At sufficiently high temperatures, both vibrational and chemical modes may lag behind the translational–rotational modes, and may also exhibit different relaxation times relative to one another, leading to a thermochemical nonequilibrium state. This delayed relaxation toward equilibrium introduces modelling challenges due to the intricate coupling between thermal and chemical processes.
% In a high-enthalpy environment, gas molecules undergo complex physicochemical processes that couple macroscopic flow properties with the microscopic internal energy modes of the gas more strongly than in classical low-enthalpy flows. These internal modes — including translation, rotation, vibration, dissociation, electronic excitation, radiation, and ionisation — absorb, store, and exchange energy at discrete, quantised levels~\cite[ch.~11]{Anderson2019}. At sufficiently high temperatures, both vibrational and chemical modes may lag behind the translational–rotational modes and may also exhibit different relaxation times relative to each other, giving rise to a thermochemical nonequilibrium state. This delayed relaxation towards equilibrium introduces additional modelling challenges due to the intricate coupling between thermal and chemical processes.
In a high-enthalpy environment, gas molecules undergo complex physicochemical processes that couple macroscopic flow properties more strongly with the microscopic internal energy modes than in classical low-enthalpy flows. These internal modes — including translation, rotation, vibration, dissociation, electronic excitation, radiation, and ionisation — absorb, store, and exchange energy at quantised levels~\cite[ch.~11]{Anderson2019}. At sufficiently high temperatures, vibrational and chemical modes can lag behind the translational–rotational modes and exhibit different relaxation times, giving rise to thermochemical nonequilibrium. This delayed relaxation introduces additional modelling challenges due to the intricate coupling between thermal and chemical processes.




In hypersonic flows, internal energy modes are predominantly endothermic, acting as heat sinks that absorb energy unlike exothermic flows such as combustion, where energy is released. The accessibility of each mode depends on the enthalpy associated with the vehicle's flight conditions~\citep{Gupta1990}, with each mode requiring a threshold energy before becoming active. As temperature increases, the internal behaviour of gas species changes significantly; a detailed depiction of this behaviour in the Earth's atmosphere is illustrated in \fref{figure:tempranges}. In conventional flows (i.e., below approximately 800~K), vibrational modes remain largely inactive, and the internal energy is governed primarily by translational and rotational contributions. Under these conditions, the gas can be treated as~\gls{CPG}, with constant specific heats and a constant ratio of specific heats, $\gamma$, more appropriately known as \textit{cold flow} in the literature. Above 800~K, vibrational excitation begins to influence thermodynamic behaviour, causing the specific heat $c_p$ and $c_v$ to vary with temperature. The gas then behaves as \emph{thermally perfect} but \emph{calorically imperfect} because $\gamma$ varies with temperature. At elevated temperatures, chemical reactions become significant, driving the dissociation and ionisation of atmospheric species over specific temperature ranges, as illustrated in \fref{figure:tempranges}~\citep{Scoggins2020}. For example, molecular oxygen (\ce{O2}) begins to dissociate at around 2500~K and is almost fully dissociated by 4000~K. These processes increase the thermochemical complexity of the flow.


\begin{figure}[h!]
    \centering
    \includegraphics[width=12cm]{resources/figures/introduction/schematic_tempranges.pdf}
    % \includegraphics[scale=0.8]{example-image-c}
    \caption{Temperature ranges at which vibrational excitation, dissociation and ionisation activate taken for air at 1-atm. Adapted from \cite{Anderson2019}. In general, the dissociation, and ionisation temperatures are directly proportional to the pressure.}
    \label{figure:tempranges}
\end{figure}


At high temperatures, additional molecular energy modes (vibrational and electronic excitation; see \fref{figure:tempranges}) necessitate a multi-temperature model to capture the relevant flow physics. Gas molecules store energy in distinct translational, rotational, vibrational, and electronic modes, each with different characteristic energy levels. These are often represented by separate temperatures: $T_t$ (translation), $T_r$ (rotation), $T_v$ (vibration), and $T_e$ (electronic excitation). A multi-temperature approach models these modes individually, providing a more accurate description of the internal energy distribution. A common simplification is to assume that the translational and rotational modes behave similarly, combining them into a single translational–rotational temperature, $T_{tr}$. When all modes are considered, the vibrational and electronic modes are often assumed to equilibrate with each other \citep{Bertin2006}, resulting in a single vibrational temperature, $T_v$. This forms the basis of Park’s classic two-temperature (2T) model~\citep{Park1990}, which remains widely used for practical high-temperature flow simulations. Although multi-temperature models can match experimental shock layer results more accurately, as shown by~\citet{Gulhan2018}, other studies suggest that detailed mode separation is only significant for expansion nozzle flows or extreme nonequilibrium conditions~\citep{ParkLee1995}. The prevailing consensus is that the increased computational cost of solving additional temperature equations is typically not justified in routine engineering practice, except in flows exhibiting pronounced nonequilibrium effects~\citep{Nagnibeda2009}.
Regardless, modelling high-enthalpy effects with sufficient fidelity is essential for predicting hypersonic flow behaviour in regimes where thermal and chemical nonequilibrium strongly influence shock-layer structure, boundary layer stability, and transition to turbulence. Capturing the correct physics and applying appropriate flow modelling strategies are therefore critical for achieving accurate and physically meaningful results.

\subsection{Flow definitions}

Chemical and vibrational processes occur over finite time scales due to molecular collisions. The relevant time scales are the vibrational relaxation time scale $\tau_{\mathrm{vib}}$, the chemical reaction time scale $\tau_{\mathrm{chem}}$, and the flow (or residence) time scale $\tau_{\mathrm{flow}}$. The relative magnitudes of these time scales give rise to distinct flow regimes, each with important implications for high-enthalpy flow behaviour. These relationships influence the rate of energy transfer between different modes, affect the onset and progression of chemical reactions, and shape the flow's overall thermal and chemical structure. The behaviour of the flow can be characterised using the Damköhler number, which compares the rate of a thermochemical process to the rate of convection. It is defined as the ratio of a characteristic chemical or vibrational time scale to the characteristic flow (or residence) time scale. For thermochemical flows, the relevant dimensionless groups are the chemical Damköhler number \gls{Dachem} and the vibrational Damköhler number \gls{Davib}, defined respectively as \citep{miromiroThesis}: \begin{equation}\label{equation:davib}
    \mathrm{Da}_{\mathrm{vib}} = \cfrac{\tau_{\mathrm{flow}}}{\tau_{\mathrm{vib}}}
\end{equation} \begin{equation}\label{equation:dachem}
    \mathrm{Da}_{\mathrm{chem}} = \cfrac{\tau_{\mathrm{flow}}}{\tau_{\mathrm{chem}}}
\end{equation} 

Depending on the relative magnitudes of the vibrational, chemical, and flow time scales, the flow regime may fall into several distinct categories. In the most general case, no assumptions are made about the relative ordering of the relevant time scales, giving rise to the \gls{TCNE} regime. All modes are of comparable magnitude, i.e., \gls{Davib} $= \mathcal{O}(1)$ and \gls{Dachem} $= \mathcal{O}(1)$. Considering the limiting case of \gls{TCE}, it is assumed that all internal modes—translational, rotational, vibrational, and chemical—adjust instantaneously relative to the flow time scales. This implies that the vibrational and chemical Damköhler numbers are much greater than one, i.e., \gls{Davib} $\gg 1$ and \gls{Dachem} $\gg 1$. For a \gls{TCF} case, vibrational and chemical modes do not have sufficient time to adjust during the flow process. As a result, the vibrational and chemical Damköhler numbers are much smaller than one, i.e., \gls{Davib} $\ll 1$ and \gls{Dachem} $\ll 1$. Isolating the Damköhler numbers allows for intermediate regimes to be examined. For example, if chemical reactions are frozen while vibrational modes adjust rapidly, then \gls{Dachem} $\ll 1$ and \gls{Davib} $\gg 1$. Conversely, if vibrational relaxation is effectively frozen while chemical reactions proceed, then \gls{Davib} $\ll 1$ and \gls{Dachem} $\gg 1$. These limiting cases provide insight into the individual roles of chemical and vibrational processes in shaping the overall flow behaviour.

\section{Boundary layer transition}
% External aerodynamic flows are inherently susceptible to environmental and surface disturbances, which can destabilise the laminar boundary layer and trigger transition to turbulence. The highly fluctuating nature of turbulent flows enhances energy and momentum mixing compared to laminar flows, resulting in steeper gradients. These large energy gradients ultimately lead to significantly higher surface heat fluxes — up to eight times greater than in a fully laminar regime~\citep{VanDriest1952}. This transition is especially critical in high-speed boundary layers, which are more prone to additional disturbances that can amplify instabilities and promote earlier breakdown. The strong shockwave compressions mean that stagnation temperatures are higher, which leads to the requirement of a \gls{TPS} unit. \gls{TPS} sizing is therefore paramount to predict the transition onset and turbulent breakdown locations on the vehicle surface for the various flow regimes encountered during the mission. Consequently, the use of inaccurate transition-modelling tools leads to designs based on the worst-case scenario, with an excessive \gls{TPS} weight allocation.

% A fundamental framework for understanding boundary-layer transition is Morkovin’s roadmap~\citep{Morkovin1994}, which describes how environmental and surface disturbances enter the boundary layer through receptivity. These disturbances generate eigenmode instabilities that amplify linearly before evolving into nonlinear interactions and secondary instabilities, ultimately leading to breakdown to turbulence. This sequence defines the pathway from laminar flow to fully developed turbulence under a given disturbance environment. A detailed analysis of these instability modes and their amplification behaviour is provided in \Cref{chapter:disturbance_growth}. Crucially, the dominant instability mechanism differs between low-speed and high-speed boundary layers, and accurately modelling these modes is essential for predicting transition onset and the eventual breakdown to turbulence. The path to full breakdown strongly depends on the nature and amplitude of the incoming disturbances: small-amplitude perturbations, such as acoustic noise or surface vibrations, typically follow the natural transition route, whereas larger surface protrusions or significant roughness can trigger bypass mechanisms. The present study focuses on isolating the natural transition path driven by small-amplitude disturbances to clarify the fundamental processes that govern instability growth and breakdown in high-speed boundary layers.
% The transition from a laminar to a turbulent boundary layer is a critical aerodynamic phenomenon with significant implications for skin friction, heat transfer, and aerodynamic drag. The destabilisation of an otherwise laminar flow into a chaotic, turbulent regime had attracted interest since the late 1800s when Osborne Reynolds hypothesised that laminar flow could become unstable under certain conditions. Over the decades, extensive research has sought to understand, predict, and control this process through theoretical and experimental efforts. Within this broader context, \citet{Morkovin1994} laid out a conceptual foundation that synthesised these developments into the modern view of the multiple paths to transition that can occur in boundary layer flows. 

% The transition from a laminar to a turbulent boundary layer is a critical aerodynamic phenomenon with significant implications for skin friction, heat transfer, and aerodynamic drag. The destabilisation of an otherwise laminar flow into a chaotic, turbulent regime has attracted interest since the late 1800s, when Osborne Reynolds hypothesised that laminar flow could become unstable under certain conditions. Over the following decades, extensive research has sought to understand, predict, and control this process through both theoretical and experimental investigations. Within this broader context, \citet{Morkovin1994} provided a conceptual framework that synthesised these developments into the modern understanding of the multiple transition paths that may arise in boundary layer flows. \fref{figure:transitionmarkovin}, reproduced from \citet{Fedorov2011}, presents a schematic overview of these paths to transition, which applies to both low-speed and high-speed boundary layers, although the specific instability modes involved differ between the two regimes.
\subsection{Paths to transition}

The transition from a laminar to a turbulent boundary layer is a critical aerodynamic phenomenon with significant implications for skin friction, heat transfer, and aerodynamic drag. Extensive research has sought to understand, predict, and control this process through both theoretical and experimental investigations. Within this broader context, \citet{Morkovin1994} provided a conceptual framework that synthesised decades of developments into the modern understanding of the multiple transition paths that can arise in boundary layer flows. \fref{figure:transitionmarkovin}, reproduced from \citet{Fedorov2011}, presents a schematic overview of these paths to transition, which is applicable to both low-speed and high-speed boundary layers, although the specific instability mechanisms differ between the two regimes.


\begin{figure}[hbt!]
    \centering
    \fontsize{9}{10}\selectfont\includesvg[width=0.50\textwidth]{resources/figures/introduction/schematic_fedorov.svg}
    \caption{\cite{Morkovin1994} roadmap to transition, diagram of pathways by which a boundary layer might transition to turbulence \citep{Fedorov2011}.}
    \label{figure:transitionmarkovin}
\end{figure}


% Transition begins with the appearance and amplification of disturbances from environmental sources such as freestream turbulence, acoustic noise, surface roughness, or shock interactions. These interact with the boundary layer, generating oscillatory modes whose amplitude, frequency, and phase define the initial conditions for instability waves that grow and ultimately breakdown. This \textit{receptivity process} governs how external disturbances are converted into instability waves, strongly influencing when and where transition occurs \citep{Morkovin1994}. In hypersonic boundary layers, key destabilising factors include freestream perturbations (acoustic, entropy, and vorticity fluctuations), surface roughness that introduces localised bypass routes, shock interactions that create unsteady pressure and velocity fields, and high-temperature effects like vibrational excitation and dissociation, all of which can significantly enhance instability growth \citep{vandereynde2015,Bitter2015}. Because these environmental disturbances are inherently complex and broadband, the receptivity problem remains largely unresolved for compressible boundary layers, despite extensive research \citep{dCerminara2017} Ma Zhong.

% A foundational framework for understanding boundary-layer transition is Morkovin’s roadmap~\citep{Morkovin1994}, which describes how environmental and surface disturbances enter the boundary layer through receptivity and generate eigenmode instabilities that amplify linearly before evolving into nonlinear interactions and secondary instabilities. This sequence defines the general pathway from laminar flow to fully developed turbulence for a given disturbance environment. Crucially, the dominant instability mechanisms differ between low-speed and high-speed boundary layers, making accurate modelling of these processes essential for predicting transition onset and breakdown. The path to turbulence strongly depends on the nature and amplitude of incoming disturbances: small-amplitude perturbations, such as acoustic noise or surface vibrations, typically follow the natural transition route, while larger roughness elements or significant freestream turbulence can trigger bypass mechanisms.
% Transition begins with the appearance and amplification of disturbances from environmental sources such as freestream turbulence, acoustic noise, surface roughness, or shock interactions. These interact with the boundary layer, generating oscillatory modes whose amplitude, frequency, and phase define the initial conditions for instability waves that grow and ultimately break down to turbulence. This conversion process, known as receptivity, governs how external disturbances generate eigenmode instabilities that amplify linearly before evolving into nonlinear interactions and secondary instabilities — a sequence formalised in Morkovin’s roadmap~\citep{Morkovin1994}. 

Transition begins with the appearance and amplification of disturbances originating from environmental sources such as freestream turbulence, acoustic noise, surface roughness, or shock interactions. These external influences interact with the boundary layer and generate oscillatory modes whose amplitude, frequency, and phase define the initial conditions for instability waves that grow and eventually break down into turbulence. This conversion process, known as receptivity, governs how external disturbances give rise to eigenmode instabilities, which amplify linearly before undergoing nonlinear interactions and secondary instabilities—a sequence formalised in Morkovin’s transition roadmap~\citep{Morkovin1994}.

% In hypersonic boundary layers, key destabilising factors include freestream perturbations (acoustic, entropy, and vorticity fluctuations), localised surface roughness, shock-wave interactions, and high-temperature effects like vibrational excitation and dissociation, all of which can significantly enhance instability growth~\citep{vandereynde2015,Bitter2015}. Because these disturbances are inherently broadband and complex, the receptivity problem remains largely unresolved for compressible boundary layers, despite extensive research~\citep{Cerminara2019}. Crucially, the dominant instability mechanisms differ between low-speed and high-speed boundary layers, making accurate modelling essential for predicting when and where transition occurs \cite{Zhong2012}. The path to turbulence strongly depends on the disturbance amplitude: small-amplitude perturbations typically follow the natural transition route, while larger roughness elements or significant freestream turbulence can trigger bypass mechanisms.




% A foundational framework for understanding boundary-layer transition is Morkovin’s roadmap~\citep{Morkovin1994}, which describes how environmental and surface disturbances enter the boundary layer through receptivity and generate eigenmode instabilities that amplify linearly before evolving into nonlinear interactions and secondary instabilities. This sequence defines the general pathway from laminar flow to fully developed turbulence for a given disturbance environment. Crucially, the dominant instability mechanisms differ between low-speed and high-speed boundary layers, making accurate modelling of these processes essential for predicting transition onset and breakdown. The path to turbulence strongly depends on the nature and amplitude of incoming disturbances: small-amplitude perturbations, such as acoustic noise or surface vibrations, typically follow the natural transition route, while larger roughness elements or significant freestream turbulence can trigger bypass mechanisms.

% The different pathways to turbulence are illustrated schematically in \fref{figure:transitionmarkovin}. In the most general case, low-amplitude oscillatory modes are triggered within the boundary layer and amplify through linear modal growth, followed by parametric instabilities and mode interactions that transfer energy between modes and lead to breakdown to turbulence (\textit{path A}). For higher disturbance levels, the superposition of perturbations may first exhibit algebraic transient growth before feeding into modal amplification (\textit{path B}). As the disturbance amplitude increases further, transient growth can lead directly to parametric instabilities and mode interactions or, if sufficiently strong, bypass these stages entirely and transition straight to breakdown (\textit{paths C} and \textit{D}). Only \textit{paths A, B,} and \textit{C} are relevant for external hypersonic boundary layers, where disturbances grow through receptivity, transient growth, and modal amplification. In contrast, \textit{path D} mainly applies to internal flows with elevated turbulence levels, while \textit{path E} represents extreme forcing such as severe surface roughness or intense freestream turbulence where linear growth mechanisms are completely bypassed. 

% Over the past decades, these physical insights have driven the development of various transition prediction methods — from empirical correlations to physics-based linear stability theory and modern numerical approaches — each aiming to capture the complex pathways by which instabilities grow and interact. The present study focuses specifically on isolating the natural transition route (\textit{path A}), driven by small-amplitude disturbances, to clarify the fundamental mechanisms that govern instability growth and breakdown in high-speed boundary layers.

% The different pathways to turbulence are illustrated schematically in \fref{figure:transitionmarkovin}. In the most general case, low-amplitude oscillatory modes are triggered within the boundary layer and amplify through linear modal growth, followed by parametric instabilities and mode interactions that transfer energy between modes and lead to breakdown to turbulence (\textit{path A}). At higher disturbance levels, the superposition of perturbations may initially exhibit algebraic transient growth before feeding into modal amplification (\textit{path B}). As the disturbance amplitude increases further, transient growth can directly trigger parametric instabilities and mode interactions or, if sufficiently strong, bypass these stages entirely and transition straight to breakdown (\textit{paths C} and \textit{D}). Among these, only \textit{paths A, B,} and \textit{C} are typically relevant for external hypersonic boundary layers, where disturbances evolve through receptivity, transient growth, and modal amplification \citep{Lefieux2021}. In contrast, \textit{path D} is more representative of internal flows with elevated turbulence levels, while \textit{path E} corresponds to extreme forcing conditions such as intense freestream turbulence or severe surface roughness where linear mechanisms are completely bypassed \cite{Reshotko2001}. These physical insights have informed the development of a range of transition prediction and control methods, from empirical correlations to physics-based linear stability theory and modern numerical simulations, each aiming to capture the complex processes through which instabilities grow and interact. The present study focuses specifically on isolating the natural transition route (\textit{path A}), driven by small-amplitude disturbances, to clarify the fundamental mechanisms governing instability growth and breakdown in high-speed boundary layers.

The different pathways to turbulence are illustrated schematically in \fref{figure:transitionmarkovin}. In the most general case, low-amplitude oscillatory modes are triggered within the boundary layer and amplify through linear modal growth, followed by parametric instabilities and mode interactions that transfer energy between modes and lead to breakdown to turbulence (\textit{path A}). At higher disturbance levels, the superposition of perturbations may initially exhibit algebraic transient growth before feeding into modal amplification (\textit{path B}). As the disturbance amplitude increases further, transient growth can directly trigger parametric instabilities and mode interactions or, if sufficiently strong, bypass these stages entirely and transition straight to breakdown (\textit{paths C} and \textit{D}). Among these, only \textit{paths A}, \textit{B}, and \textit{C} are typically relevant for external hypersonic boundary layers, where disturbances evolve through receptivity, transient growth, and modal amplification~\citep{Lefieux2021}. In contrast, \textit{path D} is more representative of internal flows with elevated turbulence levels, while \textit{path E} corresponds to extreme forcing conditions—such as intense freestream turbulence or severe surface roughness—where linear mechanisms are completely bypassed~\citep{Reshotko2001}. These physical insights have informed the development of a range of transition prediction and control methods, from empirical correlations to physics-based linear stability theory and modern numerical simulations, each aiming to capture the complex processes through which instabilities grow and interact. 

% The present study focuses specifically on isolating the natural transition route (\textit{path A}), driven by small-amplitude disturbances, to clarify the fundamental mechanisms governing instability growth and breakdown in high-speed boundary layers.

% In hypersonic boundary layers, key destabilising factors include freestream perturbations (acoustic, entropy, and vorticity fluctuations), localised surface roughness, shock-wave interactions, and high-temperature effects such as vibrational excitation and dissociation, all of which can significantly enhance instability growth~\citep{vandereynde2015,Bitter2015}. The path to transition is strongly influenced by disturbance amplitude and their spectral characteristics, and therefore understanding the process through which instabilities enter the boundary layer through receptivity is of a considerable importance. Due to the broadband and complex nature of these disturbances, the receptivity problem remains largely unresolved in compressible boundary layers, despite extensive investigation~\citep{Cerminara2019}. Importantly, the dominant instability mechanisms differ between low-speed and high-speed regimes, making accurate modelling essential for predicting when and where transition occurs~\citep{Zhong2012}. 

% In hypersonic boundary layers, key destabilising factors include freestream perturbations (acoustic, entropy, and vorticity fluctuations), localised surface roughness, shock-wave interactions, and high-temperature effects such as vibrational excitation and dissociation, all of which can significantly enhance instability growth~\citep{vandereynde2015,Bitter2015}. The path to transition is strongly influenced by both the amplitude and spectral content of these disturbances, making receptivity—the process by which external perturbations generate internal instability modes a critical area of study. Due to the broadband and complex nature of environmental inputs, the receptivity problem remains largely unresolved in compressible boundary layers, despite extensive research~\citep{Cerminara2019}. Despite this, several attempts have been made to understand the role of receptivity of free-stream disturbances in the disturbance growth, predominantly through \gls{DNS}.

% In hypersonic boundary layers, key destabilising factors include freestream perturbations (acoustic, entropy, and vorticity fluctuations), localised surface roughness, shock-wave interactions, and high-temperature effects such as vibrational excitation and dissociation, all of which can significantly enhance instability growth~\citep{vandereynde2015,Bitter2015}. The path to transition is strongly influenced by both the amplitude and spectral content of these disturbances, making receptivity, the process by which external perturbations generate internal instability modes a critical area of study. Due to the broadband and complex nature of environmental inputs, the receptivity problem remains largely unresolved in compressible boundary layers, despite extensive research~\citep{Cerminara2019}. Nonetheless, several studies have sought to elucidate the receptivity of freestream disturbances, particularly to better understand their influence on subsequent instability growth. Experimental investigation of receptivity in compressible boundary layers remains challenging, primarily due to the stringent requirement for ultra-quiet wind tunnel environments~\citep{fedorov2011transition}. The theoretical framework for such flows has received considerable attention, with studies by~\citet{Fedorov2001,Fedorov2002,MaZhong2003a,MaZhong2003b} exploring receptivity mechanisms in cold hypersonic flow over a flat plate, with \cite{Fedorov2002} accessing receptivity of wall induced disturbances on various qvectors. Further studies of broadband receptivity on flared cones were performed by \citep{Balakumar2018,HeZhong2022} concluding the role of the disturbances porpagation through a shock have altered characteristics. Recently, \citet{Varma2025} investigated the receptivity of acoustic disturbances on flared cones under thermochemical nonequilibrium conditions, concluding that receptivity can significantly influence the growth of instability modes and, consequently, the transition onset location. Nevertheless, the body of literature addressing the full transition process in high-speed flows remains limited. As highlighted by \citet{Zhong2012}, the dominant instability mechanisms differ substantially between low-speed and high-speed regimes, underscoring the need for accurate physical modelling to reliably predict the onset and location of transition.

In hypersonic boundary layers, key destabilising factors include freestream perturbations (acoustic, entropy, and vorticity fluctuations), localised surface roughness, shock-wave interactions, and high-temperature effects such as vibrational excitation and dissociation, all of which can significantly enhance instability growth~\citep{vandereynde2015,Bitter2015}. The path to transition is strongly influenced by both the amplitude and spectral content of these disturbances, making receptivity—the process by which external perturbations generate internal instability modes—a critical area of study. Due to the broadband and complex nature of environmental inputs, the receptivity problem remains largely unresolved in compressible boundary layers despite extensive research~\citep{Cerminara2019}. Experimental studies are particularly challenging owing to the stringent requirement for ultra-quiet wind tunnel environments~\citep{Fedorov2011}. As a result, theoretical frameworks have been widely explored, with~\citet{Fedorov2001,Fedorov2002,MaZhong2003a,MaZhong2003b} analysing receptivity mechanisms in cold hypersonic flow over a flat plate; \citet{Fedorov2002} in particular investigated receptivity to wall-induced disturbances across various wave-vectors. Further studies on broadband receptivity over flared cones~\citep{Balakumar2018,HeZhong2022} showed that disturbances modified by shock interactions exhibit altered characteristics. Recently, \citet{Varma2025} examined receptivity to acoustic disturbances under thermochemical nonequilibrium and found it significantly impacts instability growth and transition onset location. Nonetheless, the concept of receptivity remains widely under explored, and as noted by \citet{Zhong2012}, the dominant instability mechanisms differ between low- and high-speed regimes, underscoring the need for accurate physical modelling to reliably predict transition.



% Recetly, \cite{Varma2025} investigated the receptivity of acoustic disturbances on flared cones in thermochemical nonequilibrium, concluding that receptivity has significant effects on the growth of instabily modes and therefore, the transition onset location 

% Nevertheless, literature on the transition process remains small. As highlighted by \citet{Zhong2012}, the dominant instability mechanism differ between low-speed and high-speed regimes, outlining the need for accurate physical modelling to predict the onset location of transition. 

% Although receptivity initiates the transition process, its modelling is beyond the scope of the present study, which instead assumes the presence of small-amplitude disturbances and focuses on their subsequent amplification and breakdown.


% the dominant instability mechanisms differ markedly between low-speed and high-speed regimes, underscoring the need for accurate physical modelling to predict the onset and location of transition.


% In hypersonic boundary layers, key destabilising factors include freestream perturbations (acoustic, entropy, and vorticity fluctuations), localised surface roughness, shock-wave interactions, and high-temperature effects such as vibrational excitation and dissociation, all of which can significantly enhance instability growth~\citep{vandereynde2015,Bitter2015}. The path to transition is strongly influenced by both the amplitude and spectral content of these disturbances, making the understanding of receptivity—the process by which external perturbations generate internal instability modes—critically important. Due to the broadband and complex nature of the environmental inputs, the receptivity problem remains largely unresolved


% The path to turbulence is strongly influenced by disturbance amplitude and their spectral characteristics 


% : small-amplitude perturbations generally follow the natural transition route, whereas larger roughness elements or elevated freestream turbulence levels can trigger bypass mechanisms.




% natural transition is name for incomrpessible things, it varies for compressible flows
% The transition process has been a subject of interest since the late 19\textsuperscript{th} century, beginning with the work of Lord Rayleigh~\citep{Rayleigh1878,Rayleigh1880}, who postulated the existence of sinusoidal disturbances within the incompressible boundary layer prior to transition. This mathematical framework was extended by \citet{Tollmien1936}, who showed that incompressible flows are stable to inviscid, temporally growing disturbances unless the two-dimensional mean velocity profile contains an inflection point, $\partial_y (\partial_y U) = 0$, which provides the necessary shear layer for inviscid instabilities to extract energy from the mean flow. This theoretical prediction was confirmed experimentally by \citet{Schubauer1948}, who provided clear evidence of downstream disturbance amplification using hot-wire measurements. When the effect of viscosity was considered, these unstable waves form the classic Tollmien–Schlichting waves~\citep{Schlichting2000}. Building on this, \citet{LeesLin1946} were the first to investigate instability propagation in compressible boundary layers, although they considered only a subset of possible modes. They demonstrated theoretically that a two-dimensional, inviscid compressible boundary layer remains stable unless the mean velocity and density profiles contain a generalised inflection point, given by the condition $\partial_y (\rho, \partial_y U) = 0$, thereby extending the classic inviscid stability criterion of \citet{Tollmien1936}. They further distinguished between subsonic ($U_e - c_r < a_e$), sonic ($U_e - c_r = a_e$), and supersonic ($U_e - c_r > a_e$) disturbances, where $c_r$ is the phase speed of the disturbance, and $U_e$ and $a_e$ are the streamwise velocity and speed of sound at the boundary layer edge, respectively. This theoretical prediction was later supported by experimental observations from \cite{Laufer1960}, who helped constrain the stability boundaries for compressible flows. 

\subsection{Modal growth}
The transition process has been a subject of interest since the late 19\textsuperscript{th} century, beginning with the work of Lord Rayleigh~\citep{Rayleigh1878,Rayleigh1880}, who postulated the existence of sinusoidal disturbances within the incompressible boundary layer prior to transition. This mathematical framework was extended by \citet{Tollmien1936}, who showed that incompressible flows are stable to inviscid, temporally growing disturbances unless the two-dimensional mean velocity profile contains an inflection point, $\partial_y (\partial_y U) = 0$, which provides the necessary shear layer for inviscid instabilities to extract energy from the mean flow. This theoretical prediction was confirmed experimentally by \citet{Schubauer1948}, who provided clear evidence of downstream disturbance amplification using hot-wire measurements. When the effect of viscosity was considered, these unstable waves formed the classic Tollmien–Schlichting waves~\citep{Schlichting2000}. Building on this, \citet{LeesLin1946} were the first to investigate instability propagation in compressible boundary layers, although they considered only a subset of possible modes. They demonstrated theoretically that a two-dimensional, inviscid compressible boundary layer remains stable unless the mean velocity and density profiles contain a generalised inflection point, given by the condition $\partial_y ( \rho \partial_y U) = 0$, thereby extending the classic inviscid stability criterion of \citet{Tollmien1936}. They further distinguished between subsonic ($U_e - c_r < a_e$), sonic ($U_e - c_r = a_e$), and supersonic ($U_e - c_r > a_e$) disturbances, where $c_r$ is the phase speed of the disturbance, and $U_e$ and $a_e$ are the streamwise velocity and speed of sound at the boundary layer edge, respectively. This theoretical prediction was later supported by experimental observations from \citet{Laufer1960}, who helped constrain the stability boundaries for compressible flows. They showed that adiabatic compressible boundary layer is more stable than the incompressible counterpart. 



% Rigirous mathematical methods were first applied by \citet{LeesReshotko1962} to study two-dimensional subsonic disturbances, marking a departure from earlier asymptotic approaches. They showed that while the transition process remains structurally similar to incompressible flow, the emergence of multiple instability loops and the increasing importance of viscous dissipation terms at high Mach numbers point to fundamentally different dominant mechanisms in supersonic regimes. More comprehensive calculations were provided by \citet{Mack1963,Mack1964}, who clearly distinguished between different instability modes: the compressible Tollmien–Schlichting waves originally observed by \citet{LeesLin1946} were quantified as the first mode, while additional acoustic modes were found to emerge at higher speeds. Due to the hyperbolic nature of the governing equations in supersonic and hypersonic regimes, an infinite set of eigenvalues exists; the lowest of these, which Mack identified as the most unstable, became known as the second mode or Mack’s second mode \citep{Hudson1997}. \citet{Mack1975} demonstrated that the first mode becomes increasingly stabilised with rising Mach number, thereby rendering the second mode the dominant instability mechanism at higher supersonic and hypersonic speeds. Boundary layer instabilities may be broadly categorised as follows: at low Mach numbers, transition is primarily governed by two-dimensional Tollmien–Schlichting (T--S) waves, which are vorticity-dominated and susceptible to amplification or modal interaction in three-dimensional configurations. The first mode emerges in the transonic and early supersonic regime and is typically characterised as an oblique instability resulting from cross-stream coupling of perturbations and the attenuated upstream influence inherent to compressible flows. At higher Mach numbers, the second mode commonly referred to as the \textit{Mack mode} becomes the most unstable. This mode is acoustic in nature, with disturbances trapped between the wall and the relative sonic line. While the most amplified disturbances in two-dimensional boundary layers are themselves two-dimensional, their structure and amplification characteristics can be significantly altered in three-dimensional flows. At low Mach numbers, a two-dimensional wave is the most unstable, whereas in supersonic boundary layers, an oblique first-mode wave becomes dominant. In hypersonic boundary layers, two-dimensional second-mode waves exhibit the fastest growth. Typically, first-mode wave frequencies are in the tens of kilohertz, while second-mode waves occur at frequencies in the hundreds of kilohertz. Both linear stability theory and experimental studies have shown that wall cooling stabilises the first mode but destabilises the second mode; the second mode can even become more unstable than the first mode on a cooled wall at lower Mach numbers. Numerical simulations have also shown that the second mode can be stabilised by suction and favourable pressure gradients. 

% A rigorous mathematical framework was first developed by \citet{LeesReshotko1962} to analyse two-dimensional subsonic disturbances, marking a significant departure from earlier asymptotic approaches. Their work demonstrated that, although the transition process retains structural similarities with incompressible flows, the emergence of multiple instability loops and the increasing significance of viscous dissipation at high Mach numbers suggest fundamentally different dominant mechanisms in supersonic regimes. More comprehensive analyses were later presented by \citet{Mack1963,Mack1964}, who clearly distinguished between instability modes: the compressible equivalent of Tollmien--Schlichting waves initially observed by \citet{LeesLin1946} were identified as the first mode, while additional acoustic modes, emerging at higher speeds, were characterised as distinct instabilities. Owing to the hyperbolic nature of the governing equations in supersonic and hypersonic regimes, an infinite set of eigenvalues exists; the lowest of these, identified by Mack as the most unstable, became known as the second mode, or Mack’s second mode~\citep{Hudson1997}. \citet{Mack1975} showed that the first mode becomes increasingly stabilised with rising Mach number, thereby rendering the second mode the dominant instability mechanism at higher supersonic and hypersonic speeds. Boundary layer instabilities may be broadly categorised as follows: at low Mach numbers, transition is primarily governed by two-dimensional Tollmien--Schlichting (T--S) waves, which are vorticity-dominated and prone to amplification or modal interaction in three-dimensional configurations~\citep{miromiroThesis}. The first mode emerges in the transonic and early supersonic regime, typically as an oblique instability resulting from cross-stream perturbation coupling and the attenuated upstream influence inherent in compressible flows~\citep{Reed1996}. At higher Mach numbers, the second mode—commonly referred to as the \textit{Mack mode}—becomes dominant. This mode is acoustic in nature, with disturbances trapped between the wall and the relative sonic line. Another frequently discussed instability is the \textit{supersonic mode}, which shares similarities with the Mack mode but propagates at supersonic speeds relative to the edge velocity, radiates energy into the freestream, and exhibits a perturbation amplitude that does not decay exponentially.


% A rigorous mathematical framework was first developed by \citet{LeesReshotko1962} to analyse two-dimensional subsonic disturbances, marking a significant departure from earlier asymptotic approaches. Their work demonstrated that, although the transition process retains structural similarities with incompressible flows, the emergence of multiple instability loops and the increasing significance of viscous dissipation at high Mach numbers suggest fundamentally different dominant mechanisms in supersonic regimes. A comprehensive \gls{LST} analysis was later carried out by \citet{Mack1963,Mack1964}, providing a clear distinction between the various instability modes: the compressible equivalent of Tollmien--Schlichting waves initially observed by \citet{LeesLin1946} were identified as the first mode, while additional acoustic modes emerging at higher speeds were characterised as distinct instabilities. Owing to the hyperbolic nature of the governing equations in supersonic and hypersonic regimes, an infinite set of eigenvalues exists; the lowest of these, identified by Mack as the most unstable, became known as the second mode, or Mack’s second mode~\citep{Hudson1997}. \citet{Mack1975} showed that the first mode becomes increasingly stabilised with rising Mach number, thereby rendering the second mode the dominant instability mechanism at higher supersonic and hypersonic speeds.

% Boundary layer instabilities may be broadly categorised as follows: At low Mach numbers, where the flow is effectively incompressible, transition is primarily governed by two-dimensional Tollmien--Schlichting waves, which are vorticity-dominated and destabilised by viscous effects, and are prone to amplification or modal interactions in three-dimensional configurations~\citep{Schubauer1948}. The first mode emerges in the transonic and early supersonic regime, typically as an three-dimensional oblique instability resulting from cross-stream perturbation \citep{Mack1984} coupling and the attenuated upstream influence inherent in compressible flows. It can be viewed as the compressible counterpart of the Tollmien--Schlichting waves~\citep{Reed1996}. At higher Mach numbers, coming into the high supersonic and hypersonic regime, acoustic instabilities emerge as the dominant mode, commonly referred to as the \textit{Mack mode} or \textit{Mack's second mode} —becomes dominant \citep{Kuehl2018}. The instability arises from trapped acoustic waves between the wall and relative sonic line above the boundary layer \cite{Bitter2015}. The acoustic disturbances are convected at subsonic speeds relative to the mean flow above the sonic line, but become supersonic in the region beneath it. In two-dimensional boundary layers, the most amplified acoustic modes are strictly two-dimensional. In contrast, three-dimensional configurations favour slightly oblique modes as the most unstable~\citep{Balakumar1991}. However, \citet{Franko2013} observed that three-dimensional Mack modes undergo delayed transition relative to two-dimensional modes in zero-pressure-gradient boundary layers, suggesting that the dynamics governing their growth and interaction remain incompletely understood. A distinct class of additional modes was identified by \cite{Mack1987} that exhibits similar features to the \textit{Mack modes} but propagates supersonically with respect to the edge velocity. These modes are commonly observed in flows over highly cooled walls and, unlike the previously discussed modes, radiate energy into the freestream, resulting in a perturbation amplitude that does not decay exponentially with wall-normal distance. Further supersonic modes have also been observed in configurations where crossflow waves CITATION and Görtler vortices CITATION become prominent due to pressure gradients. The mechanisms governing this shift in dominant instability behaviour remain an active area of research. In particular, the interaction between crossflow instabilities and high-frequency acoustic modes has been investigated as a means to distinguish between competing transition scenarios~\citep{Dong2020}. 
 
% The first mode emerges in the transonic and early supersonic regime, typically as an oblique instability resulting from cross-stream perturbation coupling and the attenuated upstream influence inherent in compressible flows~\citep{Reed1996}. At higher Mach numbers, the second mode—commonly referred to as the \textit{Mack mode}—becomes dominant. This mode is acoustic in nature, with disturbances trapped between the wall and the relative sonic line. Another frequently discussed instability is the \textit{supersonic mode}, which shares similarities with the Mack mode but propagates at supersonic speeds relative to the edge velocity, radiates energy into the freestream, and exhibits a perturbation amplitude that does not decay exponentially. 

% A rigorous mathematical framework was first developed by \citet{LeesReshotko1962} to analyse two-dimensional subsonic disturbances, marking a significant departure from earlier asymptotic approaches. Their work demonstrated that, although the transition process retains structural similarities with incompressible flows, the emergence of multiple instability loops and the increasing influence of viscous dissipation at high Mach numbers indicate fundamentally different dominant mechanisms in supersonic regimes. A comprehensive \gls{LST} analysis was later carried out by \citet{Mack1963,Mack1964}, who clearly distinguished between different instability modes. The compressible equivalent of Tollmien--Schlichting waves, initially observed by \citet{LeesLin1946}, was identified as the first mode, while additional high-frequency acoustic disturbances emerging at higher speeds were classified as distinct instabilities. Owing to the hyperbolic nature of the governing equations in supersonic and hypersonic regimes, an infinite set of eigenvalues arises; the most unstable among these, identified by Mack, became known as the second mode~\citep{Hudson1997}. \citet{Mack1975} further showed that the first mode becomes increasingly stabilised with rising Mach number, thereby rendering the second mode the dominant mechanism at high supersonic and hypersonic speeds.

A rigorous mathematical framework was first developed by \citet{LeesReshotko1962} to analyse two-dimensional subsonic disturbances, marking a significant departure from earlier asymptotic approaches. Their analysis showed that, although the transition process retains structural similarities with incompressible flows, the emergence of multiple instability loops and the increasing influence of viscous dissipation at high Mach numbers point to fundamentally different dominant mechanisms in supersonic regimes. A comprehensive \gls{LST} analysis was later performed by \citet{Mack1963,Mack1964}, who clearly distinguished between different instability modes. The compressible analogue of Tollmien--Schlichting waves, initially observed by \citet{LeesLin1946}, was identified as the first mode, while additional high-frequency acoustic disturbances emerging at higher Mach numbers were classified as distinct instabilities. Due to the hyperbolic nature of the compressible governing-equations in supersonic and hypersonic flows, an infinite set of eigenvalues arises in the stability problem; the most unstable of these, identified by Mack, became known as the second mode~\citep{Hudson1997}. \citet{Mack1975} further demonstrated that the first mode becomes increasingly stabilised with rising Mach number and decreasing wall temperature, thereby rendering the second mode the dominant instability mechanism in high-speed boundary layers.


% At low Mach numbers, where the flow is effectively incompressible, transition is primarily governed by two-dimensional Tollmien--Schlichting waves, which are vorticity-dominated, destabilised by viscous effects, and prone to amplification or modal interactions in three-dimensional configurations~\citep{Schubauer1948}. The first mode emerges in the transonic and early supersonic regime, typically as a three-dimensional oblique instability driven by cross-stream perturbation coupling and reduced upstream influence~\citep{Mack1984,Reed1996}. At higher Mach numbers, the second mode—commonly referred to as the \textit{Mack mode}—becomes dominant~\citep{Kuehl2018}. This instability arises from acoustic waves trapped between the wall and the relative sonic line above the boundary layer~\citep{Bitter2015}, with disturbances convected subsonically above the sonic line and supersonically beneath it. In two-dimensional boundary layers, the most amplified second-mode disturbances remain strictly two-dimensional, while in three-dimensional flows, the most unstable waves exhibit slight obliqueness~\citep{Balakumar1991}. However, \citet{Franko2013} observed that three-dimensional Mack modes undergo delayed transition relative to their two-dimensional counterparts in zero-pressure-gradient boundary layers, suggesting that the mechanisms governing their amplification remain incompletely understood. A distinct class of additional modes, referred to as \textit{supersonic modes}, was further identified by \citet{Mack1987}. These share some characteristics with the second mode but propagate supersonically with respect to the edge velocity. Supersonic modes are commonly observed in flows over highly cooled walls and, unlike the previously discussed instabilities, radiate energy into the freestream—resulting in a perturbation amplitude that does not decay exponentially with wall-normal distance. 

At low Mach numbers, where the flow is effectively incompressible, transition is primarily governed by two-dimensional Tollmien--Schlichting waves, which are vorticity-dominated, destabilised by viscous effects, and prone to amplification or modal interactions in three-dimensional configurations~\citep{Schubauer1948}. The first mode emerges in the transonic and early supersonic regime, typically as a three-dimensional oblique instability driven by cross-stream perturbation coupling and reduced upstream influence~\citep{Mack1984,Reed1996}. At higher Mach numbers, the second mode—commonly referred to as the \textit{Mack mode}—becomes dominant~\citep{Kuehl2018}. This instability arises from acoustic waves trapped between the wall and the relative sonic line above the boundary layer~\citep{Bitter2015}, with disturbances convected subsonically above the sonic line and supersonically beneath it. In two-dimensional boundary layers, the most amplified second-mode disturbances remain strictly two-dimensional, whereas in three-dimensional flows, the most unstable waves exhibit slight obliqueness~\citep{Balakumar1991}. However, \citet{Franko2013} observed that three-dimensional Mack modes undergo delayed transition relative to their two-dimensional counterparts in zero-pressure-gradient boundary layers, suggesting that the mechanisms governing their amplification remain incompletely understood. A distinct class of additional modes, referred to as \textit{supersonic modes}, was later identified by \citet{Mack1987}. These share some characteristics with the second mode but propagate supersonically with respect to the edge velocity. Supersonic modes are commonly observed in flows over highly cooled walls and, unlike the previously discussed instabilities, radiate energy into the freestream, resulting in a perturbation amplitude that does not decay exponentially with wall-normal distance.

Further  modes have also been identified in configurations with pressure gradients, where crossflow instabilities and Görtler vortices become prominent features of the transition process. The mechanisms governing this shift in dominant instability behaviour remain an active area of research. In particular, the interaction between crossflow modes and high-frequency acoustic disturbances has been investigated as a means of distinguishing competing transition scenarios~\citep{Dong2020}. Considering a general \gls{ZPG} flat plate boundary layer, the prevailing consensus in the literature is that the dominant instability mode is \textit{Mack's second mode}, typically associated with frequencies in the hundreds of kilohertz range. Among the various instability mechanisms, these high-frequency acoustic modes are particularly relevant in hypersonic regimes due to their rapid amplification, strong wall-normal gradients, and central role in triggering transition.

    
% Since their inception, \textit{Mack’s second modes} have been the focus of extensive numerical and experimental studies aimed at understanding the behaviour of these acoustic instabilities. \citet{Kendall1975} was the first to demonstrate the existence of Mack’s second mode experimentally. Furthermore, the work of \citet{Lysenko1984} corroborated that the second mode remains the primary instability mechanism for Mach numbers of 4 and above, while also verifying the effects of wall cooling on both the first and second modes. \citet{Erlebacher1990} performed temporal \gls{DNS} to investigate the behaviour and stability of Mack modes, extending the \gls{LST} framework of \citet{Mack1963} and confirming the existence of these additional acoustic modes. 

Since their inception, \textit{Mack’s second modes} have been the subject of extensive numerical and experimental investigations aimed at characterising the behaviour of these high-frequency acoustic instabilities. \citet{Kendall1975} was the first to confirm the existence of the second mode experimentally. Subsequently, \citet{Lysenko1984} demonstrated that the second mode remains the dominant instability mechanism for Mach numbers above 4, while also highlighting the influence of wall cooling on both first and second modes. Building on the \gls{LST} framework established by \citet{Mack1963}, \citet{Erlebacher1990} conducted temporal \gls{DNS} to examine the stability characteristics of the second mode and their subsequent breakdown to turbulence, and confirmed the presence of additional acoustic instabilities. Despite decades of research, a comprehensive understanding of Mack modes—particularly their nonlinear interactions and behaviour under thermally excited, real-gas conditions—remains elusive. The linear and early nonlinear stages, along with the role of thermochemical nonequilibrium in typical flight scenarios, therefore remain key areas of interest and will be explored in detail in later sections.

% To date, most investigations of Mack modes have been performed in cold hypersonic flows, where thermal and chemical excitations are neglected. Although recent works have begun to explore the influence of thermochemical nonequilibrium on second-mode instability, a complete physical picture has yet to emerge.


% Since their inception, \textit{Mack’s second modes} have been the focus of extensive numerical and experimental studies aimed at understanding the behaviour of these acoustic instabilities. \citet{Kendall1975} was the first to demonstrate the existence of Mack’s second mode experimentally. Furthermore, the work of \citet{Lysenko1984} corroborated that the second mode remains the primary instability mechanism for Mach numbers of 4 and above, while also verifying the effects of wall cooling on both the first and second modes. \citet{Erlebacher1990} performed temporal \gls{DNS} to investigate the behaviour and stability of Mack modes, extending the \gls{LST} framework of \citet{Mack1984} and confirming the existence of additional acoustic modes. 


% Since their inception, \textit{Mack’s second modes} have been the focus of extensive numerical and experimental studies aimed at understanding the behaviour of these acoustic instabilities. \citet{Kendall1975} was the first to demonstrate the existence of Mack’s second mode experimentally. Furthermore, the work of \citet{Lysenko1984} corroborated that the second mode remains the primary instability mechanism for Mach numbers of 4 and above, while also verifying the effects of wall cooling on both the first and second modes. \citet{Erlebacher1990} performed temporal \gls{DNS} to investigate the behaviour and stability of Mack modes, extending the \gls{LST} framework of \citet{Mack1984} and confirming the existence of additional acoustic modes. Despite decades of research, a comprehensive understanding of Mack modes, their non-linear interactions, and their behaviour under thermally excited, real-gas conditions remains elusive. This gap underscores the need for further studies that couple high-fidelity simulations with experiments to resolve these complex instability mechanisms. For boundary layer flows with high stagnation enthalpy, the temperature within the interior of the boundary layer can become sufficiently large that the vibrational energy modes of polyatomic molecules are significantly excited, rendering the calorically perfect gas assumption invalid. At even higher enthalpies, dissociation reactions may also become important. Furthermore, the timescales associated with chemical reactions or vibrational–translational energy transfer can be comparable to the flow timescales, resulting in a thermochemical nonequilibrium flow. A number of researchers have investigated these nonequilibrium effects on the stability of boundary layers, but in some cases, conflicting results have been reported. For example, \citet{Hudson1997} concluded that “thermal nonequilibrium has a destabilizing effect,” whereas \citet{Bertolotti1998} found that “vibrational relaxation has a large destabilizing influence.” By contrast, \citet{Johnson1998} determined that “translational–vibrational energy transfer can be either stabilizing or destabilizing.” Recent works of \cite{Bitter2015} and \cite{miromiro2018} have extended the understanding of the modes in the linear regime, but their interaction in nonlinear regime is an area that has not been fully explored, let alone in flows that are thermally excited and exhibit different characteristics. 

% \newpage

\subsection{Breakdown to turbulence}


Following the initial amplification of instbaility modes predicted by linear theories 

% Following the initial amplification of instability modes, breakdown to turbulence is characterised by the onset of nonlinear interactions that distort coherent flow structures and generate a hierarchy of smaller-scale vortical motions.
Following the initial amplification of instability modes predicted by linear theories—and, where relevant, transient non-modal growth \citep{ButlerFarrell1991,Trefethen1993,SchmidHenningson2001} their amplitudes eventually become sufficient to distort the laminar base flow and trigger nonlinear interactions among modes. These interactions excite additional spatial and temporal scales, inducing harmonics and subharmonics of the dominant disturbance frequencies, thereby accelerating breakdown of the laminar base flow into the chaotic motion and broadband spectrum characteristic of turbulence. The evolution and breakdown of localised three-dimensional disturbances in Poiseuille flow were first mapped in detail by \citet{Henningson1993}, showing that small disturbances undergo substantial transient growth by the lift-up mechanism \citep{Landahl1975}, form elongated streamwise streaks, and then evolve nonlinearly to breakdown. The streak-induced, spanwise-modulated shear layers that arise at finite amplitude support secondary instabilities, which couple the primary content with harmonics and subharmonics, proliferate scales, and precipitate transition \citep{Herbert1988,Kachanov1994,ReedSaricArnal1996}. 



% Following the initial amplification of instability modes predicted by linear theories and, where relevant, transient non-modal growth
Following the initial amplification of instability modes predicted by linear theories, relevant transient non-modal growth~\citep{ButlerFarrell1991} their amplitudes eventually become sufficient to distort the laminar base flow and trigger nonlinear interactions among modes. These interactions excite adiditional spatial and temporal scales, including harmonics and subharmonics of the dominant disturbance frequencies, thereby accelerating breakdown 

References:
\citep{ButlerFarrell1991}, 
\citep{Trefethen1993}
\citep{SchmidHenningson2001}
\citep{Hanifi1996}
\citep{Andersson2001}
\citep{BrandtHenningson2002}
\citep{alam_sandham_2000}


% Following the initial amplification of instability modes predicted by linear theories, their amplitudes eventually becomes sufficient to distort the laminar base flow and trigger nonlinear interactions among modes. These interactions excite additional spatial and temporal scales inducing harmonics and subharmonics of the dominant disturbance frequencies thereby accelerating breakdwon of laminar  base flow into the chaotic motion and broadband spectrum characteristic of turbulence.  The evolution and breakdown of localised three-dimensional disturbances in Poiseuille flow were first mapped in detail by \citet{Henningson1993} showing that small disturbances undergo substantial transient growth by the lift-up mechanism \citep{Landahl1975}, form elongated streamwise streaks, and then evolve non-linearly to breakdown. 

\cite{Hanifi1996} extended the previous frameworks to compressible regime, looking into temporal transient growth analyses. 

. The finite amplitude localised disturbances were shown to undergo transition through the lift-up phenomenon \citep{Landahl1975}, 

% \cite{Henningson1993} were the first to study the spatial evolution of localised disturbances, their subsequent breakdown in three-dimensional Poiseuille flow.


The evolution of a three-dimensional localized disturbance in channel flow was first examined by Henningson, Lundbladh \& Johansson (1993), who showed that a small-amplitude disturbance undergoes linear transient growth via the lift-up mechanism \citep{Landahl1975}, whereby the disturbance three-dimensionality generates energetic streamwise streaks, amplifies wall-normal perturbation vorticity, and ultimately preconditions the flow for secondary instability and breakdown.




% Path E in \fref{figure:transitionmarkovin} denotes disturbances of sufficiently large initial amplitude that they enter the nonlinear regime without prior exponential amplification, a route commonly referred to as bypass transition.



As the disturbances amplify according to the above mentioned mechanisms, their amplitude eventually becomes large enough to modufy the laminar flow profile and cause non-linear interactions between modes. Because of these nonlinear interactions, additional legnth and time scales are excited at harmonics of the dominant disturbance frequencies and the flow rapidly breaks down into the chaotic motion and broad-band frequency spectrum of turbulence. Path E in \fref{figure:transitionmarkovin} accounts for disturbances that are initially of such large amplitude that no amplification is needed in order to reach the onlinear regime, which is often called bypass transition. Although initial efforts to predict transition to turbulence focused on the modal instability mechanisms, it was soon realised that modal instabilities could not explain all the phenomenon observed in experiments. In particular, certain canonical flows, such as pipe flow and Couette flow, have no unstable modes, and yet transition to turbulence is readily achieved in experiments. Other flows, like plane Poiseuille flow, transition at low Reynolds numbers for which the linear theory predicts that the flow is stable. 

The modal analysis 

\cite{Herbert1988} talks about incompressible secondary instabilities and breakdown scenarios that include subharmionic and fundamental secondary instabilities. 

\cite{SchmidHenningson2001} talks about transition in three different transition mechanisms for incompressible and supersonic flow. Oblique breakdown starts with the interaction of two oblique waves (either first mode or second mode), which interacts non-linearly to produce a dominant spanwise varying stationary mode. This mode contains large streamwise vorticity and leads to velocity and temperature streaks in the boundary layer. Higher modes are also generated nonliearly and the boundary layer completely transitions to turbulence. Subharmonic instability begins with the growth of a 2D wave, eitehr a first mode or a second mode instability. Onbce the 2D wave is large enough, an obliquye subharmonic mode ( frequency half of primary mode and with a spanwise wavenumber) grows rapidly. The subharmonic modes then ineract to produce a stationary mode ( similar to oblique transition). Fundamental transition is similar to
subharmonic transition, except a three dimensional disturbance with the same frequency as the primary disturbance is generated. 


In a boundary layer flow, after the initial transient stage is covered by modal growth. The precursor of transition to turbulence is given by the secondary instability or nonmodal growth that breaakdowns to turbulence.  The transition process over a flat plate boundary alyer is shown in the figure below. This scenario was investigated almost simultaneously in the theoretical work of Herbert  the experimetnal work of Kachavov and Levchenko and DNS works of Kleiser Schumann, Henningson. When the most unstable of primary instability grows in amplitude, for an incomrpessible flow, this is tollmein schiting waves need to reach 1\% of the  frestream. This new baseflow can be unstable to exponentially small disturbances, particularly three-dimensional waves in the boundary layer Herbert 1988. Since the new formulation will involved equations with periodic coefficients, Floquet theory analyses this newly modified state. Physically, secondary instabilities represent lambda-shaped voritces formed during the late stages of transiiton near the boundary layer edge. In this matter, three clasees of linear secondary instabilities have been idnetified experimentally by Kachanov and Levchenko and Saric and Thomas. 


The first fundamental instability/resonance, aslo referred to as K-type instbailiity due to the experimental work of Klebanoff. In this case, the secondary waves have the exact streamwise wavenumber as the primary waves and are aligned with each other. The second one is subharmonic resonance also known as H-type instabilitye ( Herbert 1988). In this case, the streamwise wavenumber of the secondary instability is twice the primary and forms a staggered vortices pattern. DNS simulations performed by Spalart and Yang 1987 showed that H type instbaility is the most unstable secondary instability in an incompressible boundary alyer. However, K-type has been found to be dominant in experiments. This explained due to low amplitudes streamwise voriticity in the background flow. The third class is referred to as O-type instbaility (oblique breakdown) and was discovered by Thummm, Fasel. This breakdown starts from the nonlinear interaction of two oblique instability waves with equal wave angles but in opposite directions. 

Now, let's chat about the compressible trnasition mechanism, there will be the extra things with Y-shaped transition process and it was studies in the past by Sandham and Adamas, Sandham et al, Adamns and Kleiser, Pruett and Zhang, Mayer - oblique breakdown. For maintly cold  wall flows, recently studies of 

Due to the difficulties of experimental studies of hypersonic boundary layer transition, DNS can give insight into the potential mechanisms for breakdown to turbulence and the overshoot in wall heat transfer. This insight can then be used to guide further exploration in both computations and experiments



% \cite{Mack1975} further showed that the first mode becomes increasingly stabilised with rising Mach number, thereby rending the second-mode the dominant instability mechanism at higher supersonic and hypersonic speeds. Boundary layer instabilities may be broadly categorised as follows: at low Mach numbers, transition is primarily governed by the two-dimensional Tollmien--Schlichting waves, which are vorticity dominated and prone to amplification or modal interaction in three-dimensional configuration \citep{miromiroThesis}. The first mode emerges as the dominant mode in the transonic and early supersonic regime, typically as an oblique instability resulting from cross-stream perturbation coupling and the attenuated upstream influence inherent in compressible flows \citep{Reed1996}. \citet{Mack1975} further showed that the first mode becomes increasingly stabilised with rising Mach number, thereby rendering the second mode the dominant instability mechanism at higher supersonic and hypersonic speeds. Boundary layer instabilities may be broadly categorised as follows: at low Mach numbers, transition is primarily governed by two-dimensional Tollmien--Schlichting (T--S) waves, which are vorticity-dominated and prone to amplification or modal interaction in three-dimensional configurations. The first mode emerges in the transonic and early supersonic regime, typically as an oblique instability resulting from cross-stream perturbation coupling and the attenuated upstream influence inherent in compressible flows. At higher Mach numbers, the second mode---commonly referred to as the \textit{Mack mode}---becomes dominant. This mode is acoustic in nature, with disturbances trapped between the wall and the relative sonic line. Another commonly discussed instability is the \textit{supersonic mode}, which shares notable similarities with the Mack mode; however, it propagates at supersonic speeds relative to the edge velocity, radiates energy into the freestream, and exhibits a perturbation amplitude that does not decay exponentially. 








% As will be shown in Chapter 5, vibrational relaxation can introduce two competing effects that may lead to net stabilization or destabilization, depending on the specific flow conditions. Chapter 5 investigates these mechanisms in detail and clarifies the conditions under which net stabilization or destabilization is likely to occur.

% While the most amplified disturbances in two-dimensional boundary layers are also two-dimensional, their structure and amplification characteristics can vary significantly in three-dimensional flows. At low Mach numbers, the most unstable waves remain two-dimensional, whereas in supersonic regimes, oblique first-mode waves are typically more unstable. In hypersonic flows, two-dimensional second-mode waves exhibit the fastest growth. First-mode waves generally occur at frequencies in the tens of kilohertz, while second-mode waves lie in the hundreds of kilohertz range. Both linear stability theory and experimental studies have shown that wall cooling stabilises the first mode but destabilises the second; indeed, on a cooled wall, the second mode may become more unstable than the first at lower Mach numbers. Numerical simulations have further demonstrated that the second mode can be stabilised through suction and favourable pressure gradients.



% \citet{Mack1975} further showed that the first mode is stabilised with increasing Mach number, making the second mode dominant at higher speeds. Strictly speaking, the instabilities are divided as fllows: there are the T-S waves, which are 2D vorticity waves the main reason for transition in low speed flows, when they interact with on modally in 3 dimensions. The first mode arises in transonic and early supersonic regime, an oblique mode, usually described/arises by cross-stream mutual interaction between disturbances and
% decreased upstream influence. The most unstable mode at high supersonic regime and hypersonic regime is the 2nd mode or comonly known as the \textit{Mack} mode. They are considered acoustic instabilites trapped between the wall and the relative sonic line. In 2D BL, most unstable is a purely 2D more, but their behaviour changes in 3D where ... 












% The acoustic second-mode instabilities produce a higher growth rates than the first mode Tollmien–Schlichting waves at higher Mach numbers, and are destabilised by cold walls, high Mach numbers, ... 


% \cite{Mack1969} performed a comprehensive linear stability analysis of compressible boundary layers in the high-supersonic and low-hypersonic regimes, significantly advancing the understanding of flow stability under these conditions.

% This early understanding of compressible flow stability set the stage for the more comprehensive framework developed by Mack, who formalised the linear stability theory for high-speed boundary layers and identified the distinct instability modes that arise in the hypersonic regime. A detailed framework for compressible high-speed boundary layer stability analysis can be found in \citet{Mack1984} and the references therein. Notably, \citet{Mack1969} showed that inviscid instability increases with Mach number as the generalised inflection point moves away from the wall, while viscous instability decreases and eventually disappears for two-dimensional waves at $Ma > 3$, implying that viscosity has a stabilising effect at high speeds. 

% Within this framework, Mack formally classified these modes, distinguishing the first-mode instabilities — primarily viscous in nature and equivalent to the classic Tollmien–Schlichting waves — from the second-mode instabilities, which are essentially trapped acoustic waves that become dominant at high Mach numbers.



% A detailed framework for compressible high-speed boundary layer stability analysis can be found in \citet{Mack1984} and the references therein. \cite{Mack1969} showed that inviscid instability increases with Mach number as the generalised inflection point moves away from the wall, while viscous instability decreases and eventually disappears for two-dimensional waves at $Ma > 3$, so that viscosity has a stabilising effect in high-speed flows. 

% Most of the groundwork done for  

% Ok, hear me out, most of the groundwork done for transition prediction was on incompressible flows which is not completely true but we say that anyway. The first studies about boundary layer stability was performed my Lees and Lin, who were quantifying the inflection point phenomenon, then came the Tollmein-Schliting waves, confirmed by further experiments by Schuberger or whatever their names are. It was not until Mack studies the stability of supersonic boudnary layer that we have a rigiorous definition of boundary layers and what they actually are. Mack performed stability analysis on supersonic base flows and characterised instability modes into numbers (first, second and third mode, so on). A detailed description of these instabilities is discussed in chapter 4, but the general gist of it is, first modes are promiment in subsonic incompressible flows, often associated with lower frequencies and the second mode is an acoustic trapped wave that is a feature of flows past supersonic regime. 


% The transition mechanism must be well understood for regime Zhong2012


% In the most general case, low-amplitude oscillatory modes are triggered within the boundary layer and amplify through linear modal growth, followed by parametric instabilities and mode interactions that transfer energy between modes and lead to breakdown to turbulence; this transition pathway is depicted in \fref{figure:transitionmarkovin} as \textit{path A}. For higher disturbance levels, the superposition of perturbations may first exhibit algebraic transient growth before feeding into modal amplification (\textit{path B}). As the disturbance level increases further, transient growth can lead directly to parametric instabilities and mode interactions or, if sufficiently strong, bypass these stages entirely and transition straight to breakdown (\textit{paths C} and \textit{D}). Only \textit{paths A, B,} and \textit{C} are relevant for external hypersonic boundary layers, where disturbances grow through receptivity, transient growth, and modal amplification. \textit{Path D} mainly applies to internal flows with elevated turbulence levels, while \textit{path E} represents extreme forcing such as severe surface roughness or intense freestream turbulence where linear growth mechanisms are completely bypassed.


% A foundational framework for understanding boundary-layer transition is Morkovin’s roadmap~\citep{Morkovin1994}, which describes how environmental and surface disturbances enter the boundary layer through receptivity and generate eigenmode instabilities that amplify linearly before evolving into nonlinear interactions and secondary instabilities. This sequence defines the pathway from laminar flow to fully developed turbulence for a given disturbance environment. Crucially, the dominant instability mechanisms differ between low-speed and high-speed boundary layers, making accurate modelling of these processes essential for predicting transition onset and breakdown. The path to turbulence strongly depends on the nature and amplitude of incoming disturbances: small-amplitude perturbations, such as acoustic noise or surface vibrations, typically follow the natural transition route, while larger roughness elements or protrusions can trigger bypass mechanisms. Over the past decades, these physical insights have driven the development of various transition prediction methods — from empirical correlations to more physics-based linear stability theory and modern numerical approaches — each attempting to capture the complex pathways by which instabilities grow and interact. The present study focuses on isolating the natural transition route driven by small-amplitude disturbances to clarify the fundamental mechanisms that govern instability growth and breakdown in high-speed boundary layers.


% A foundational framework for understanding boundary-layer transition is Morkovin’s roadmap~\citep{Morkovin1994}, which describes how environmental and surface disturbances enter the boundary layer through receptivity and generate eigenmode instabilities that amplify linearly before evolving into nonlinear interactions and secondary instabilities. This sequence defines the pathway from laminar flow to fully developed turbulence for a given disturbance environment. Crucially, the dominant instability mechanisms differ between low-speed and high-speed boundary layers, making accurate modelling of these processes essential for predicting transition onset and breakdown. The path to turbulence strongly depends on the nature and amplitude of incoming disturbances: small-amplitude perturbations, such as acoustic noise or surface vibrations, typically follow the natural transition route, while larger roughness elements or protrusions can trigger bypass mechanisms. Over the past decades, these physical insights have driven the development of various transition prediction methods — from empirical correlations to more physics-based linear stability theory and modern numerical approaches — each attempting to capture the complex pathways by which instabilities grow and interact. The present study focuses on isolating the natural transition route driven by small-amplitude disturbances to clarify the fundamental mechanisms that govern instability growth and breakdown in high-speed boundary layers.


% \begin{figure}[hbt!]
%     \centering
%     \fontsize{9}{10}\selectfont\includesvg[width=0.50\textwidth]{resources/figures/introduction/schematic_fedorov.svg}
%     \caption{\cite{Morkovin1994} roadmap to transition, diagram of pathways by which a boundary layer might transition to turbulence \citep{Fedorov2011}.}
%     \label{figure:transitionmarkovin}
% \end{figure}


% The natural transition process has been a subject of interest since the late 19\textsuperscript{th} century, beginning with the work of Lord Rayleigh~\citep{Rayleigh1878,Rayleigh1880}, who postulated the existence of sinusoidal disturbances within the incompressible boundary layer prior to transition. This mathematical framework was extended by \citet{Tollmien1936}, who showed that incompressible flows are stable to inviscid, temporally growing disturbances unless the two-dimensional mean velocity profile contains an inflection point. This theoretical prediction was confirmed experimentally by \citet{Schubauer1948}, who provided clear evidence of downstream disturbance amplification using hot-wire measurements. When the effect of viscosity is included, these unstable waves are known as the classic Tollmien–Schlichting (TS) waves~\citep{Schlichting2000}. Building on this, \citet{LeesLin1946} were the first to investigate instability propagation in compressible boundary layers, although they considered only a subset of possible modes. They demonstrated that a two-dimensional, inviscid compressible boundary layer remains stable unless the mean velocity and density profiles contain a generalised inflection point, analogous to the condition identified by \citet{Tollmien1936}. This insight established the well-known inflection point criterion for compressible flows.


% A detailed framework for compressible high-speed boundary layer stability analysis can be found in \citet{Mack1984} and the references therein. \cite{Mack1969} showed that inviscid instability increases with Mach number as the generalised inflection point moves away from the wall, while viscous instability decreases and eventually disappears for two-dimensional waves at $Ma > 3$, so that viscosity has a stabilising effect in high-speed flows. It was also within this work that Mack formally classified these instability modes, distinguishing the first-mode instabilities—primarily viscous in nature and corresponding to the classical Tollmien–Schlichting waves—from the second-mode instabilities, which are essentially inviscid acoustic waves trapped within the boundary layer and become dominant at high Mach numbers. Mack performed \gls{LST} on supersonic flows and found that there are infinite number of eigenvalues due to the hyperbolic nature of the equations. The lowest order is called the second mode. As the Mach number increases, the first mode maximum growth
% rate decreases. It even becomes less unstable than the second mode which becomes the
% dominant instability mode for M > 4. In low speed flows, a 2D wave is most unstable
% whereas an oblique first mode wave is most unstable in a supersonic boundary layer (M
% < 4). In a hypersonic boundary layer, the 2D second-mode waves grow faster. Frequencies of first-mode waves are around tens of kilohertz while second mode waves can be found around hundreds of kilohertz. Both LST and experiments showed that wall cooling had the effect of stabilizing the first mode and destabilizing the second mode. The second mode could even be more unstable than the first mode on a cooled wall at smaller Mach numbers. Numerical simulations showed that the second mode could be stabilized by suction and favourable pressure gradients.

% Unlike the Tollmien-Schlichting waves mentioned above which can be stabilized by
% wall cooling, the acoustic modes are destabilized, an eect which is mainly caused by
% the reduction in boundary layer thickness that occurs when the wall is cooled. This
% destabilization can be thought of in the following way: as the boundary layer becomes
% thinner, the wavelength of the second mode waves also decreases since they are reverberating within the boundary layer. The amplication per unit wavelength remains
% approximately constant, and hence the amplication per unit distance (the growth
% rate) increases as the boundary layer is made thinner, since there are more waves per
% unit distance. This explanation is demonstrated quantitatively in Section 5.4.2.

% For boundary layer fows that have a high stagnation enthalpy, the temperature
% in the interior of the boundary layer can become large enough that the vibrational
% energy of polyatomic molecules becomes signicant, such that the gas cannot
% be treated as calorically perfect. Dissociation reactions can also become signi-
% cant at suciently high enthalpies. Furthermore, the rates of chemical reactions or
% vibrational-translational energy transfer can be comparable to 
% ow timescales, leading to a nonequilibrium ow. A number of researchers have considered these nonequilib turbations appears in the 
% ow, for instance, behind a small, sharp-edged roughness
% element. The vertical velocity generated by these vortices then carries low momentum uid up from the wall and high momentum 
% uid down from the freestream, thereby producing streaks of high and low streamwise velocity. The existence of this mechanism
% has been demonstrated in experiments involving roughness elements placed in
% incompressible boundary layers. Such experiments have reproduced many features
% predicted by transient growth calculations, including the downstream evolution of
% disturbance kinetic energy and the shape of the amplied disturbances (White, 2002,
% White et al., 2005, Ergin and White, 2006).




% In the most general case, \gls{TCNE}, no assumptions are made about the ordering of the time scales which are of comparable magnitude, i.e., $\mathrm{Da}_{\mathrm{vib}} =\mathcal{O}(1)$ and $\mathrm{Da}_{\mathrm{chem}} =\mathcal{O}(1)$. 

% When vibrational relaxation is much faster than the flow time scale, and both are faster than the chemical reactions, the regime is described as chemically frozen and in \gls{TEQ}, where $\tau_{\mathrm{vib}} \ll \tau_{\mathrm{flow}} \ll \tau_{\mathrm{chem}}$. If both vibrational and chemical processes relax much more rapidly than the flow, i.e., $\tau_{\mathrm{chem}} \sim \tau_{\mathrm{vib}} \ll \tau_{\mathrm{flow}}$, the flow is in \gls{TCE}. On the other hand, when vibrational relaxation is much slower than both flow motion and chemical activity, i.e., $\tau_{\mathrm{vib}} \gg \tau_{\mathrm{flow}} \sim \tau_{\mathrm{chem}}$, the flow is in a thermal equilibrium but chemical non-equilibrium (TEQ–CNE) state. Conversely, if chemical relaxation is fast compared to flow and vibrational processes, i.e., $\tau_{\mathrm{chem}} \ll \tau_{\mathrm{flow}} \sim \tau_{\mathrm{vib}}$, the flow exhibits thermal non-equilibrium with chemical equilibrium (TNE–CE). These classifications are essential for guiding the appropriate choice of physical models in high-enthalpy flow simulations.

% Depending on the relative magnitudes of these time scales, the flow regime can be classified into the following scenarios: 

% \begin{itemize}
%     \item \textbf{Thermal and chemical non-equilibrium} TCNE is the most general case, with no assumptions about the relation between the time scales, i.e. $\tau{vib} \sim \tau{chem} \sim \tau{flow}$. 
%     \item \textbf{Chemically frozen and thermal equilibrium} CF-TEQ occurs if the vibrational relaxation time occurs substantially faster than the flow times scale and both are smaller than the chemical timescale $\tau{vib} \ll \tau{flow} \ll \tau{chem}$.
%     \item \textbf{Thermal and chemical equilibrium} TCEQ occurs when both the chemical and vibrational time scales are significantly faster than the flow time scale: $\tau{chem} \sim \tau{vib} \ll \tau{flow}$
    
%     \item \textbf{Thermal equilibrium, chemical non-equilibrium} TEQ-CNE is assumed when the chemical activity and the flow motion are significantly faster than the vibrational relaxation: $\tau{vib} \ll \tau{flow} \sim \tau{chem}$. 
    
%     % The two-temperature assumption (\eref{equation:twotemperature}) could be reduced down to a one temperature model when this condition occurs. 
%     \item \textbf{Thermal non-equilibrium, chemical equilibrium} TNE-CE, where the chemical relaxation is significantly slower than the chemical process and flow motion $\tau{chem} \ll \tau{flow} \sim \tau{vib}$, the system is assumed to be in a thermally non-equilibrium state. 
% \end{itemize}



% \section{Laminar turbulent transitions}
% % \section{Transitional flows}
% The transition from a laminar to a turbulent boundary layer is a critical aerodynamic phenomenon with significant implications for skin friction, heat transfer, and aerodynamic drag. The destabilisation of an otherwise laminar flow into a chaotic, turbulent regime had attracted interest since the late 1800s when Osborne Reynolds first demonstrated that laminar flow could become unstable under certain conditions. Over the decades, extensive research has sought to understand, predict, and control this process through theoretical and experimental efforts. Within this broader context, \citet{Morkovin1994} laid out a conceptual foundation that synthesised these developments into the modern view of the multiple paths to transition that can occur in boundary layer flows. \fref{figure:transitionmarkovin}, reproduced from \citet{Fedorov2011}, presents a schematic overview of these paths to transition, which applies to both low-speed and high-speed boundary layers, although the specific instability modes involved differ between the two regimes.


% \begin{figure}[t!]
%     \centering
%     \fontsize{9}{10}\selectfont\includesvg[width=0.50\textwidth]{resources/figures/introduction/schematic_fedorov.svg}
%     \caption{Diagram of pathways by which a boundary layer might transition to turbulence \citep{Fedorov2011}.}
%     \label{figure:transitionmarkovin}
% \end{figure}

% % The transition begins with the appearance and subsequent amplification of disturbances, which can arise from environmental sources such as freestream turbulence, acoustic noise, surface roughness, or shock interactions, depending on the flow regime. These environmental disturbances interact with the boundary layer, appearing as fluctuations in the base flow that generate oscillatory modes whose amplitude, frequency, and phase set the initial conditions for the instability waves that develop and ultimately lead to breakdown. This process, known as the receptivity mechanism, converts external disturbances into instability waves; their initial amplitude plays a key role in determining how and when the transition occurs. Because environmental disturbances are often complex and contain a wide range of information, the receptivity problem remains largely unsolved for compressible boundary layers, although several studies have tried to tackle the issue at hand. 

% % Transition begins with the appearance and subsequent amplification of disturbances originating from environmental sources such as freestream turbulence, acoustic noise, surface roughness, or shock interactions, depending on the flow regime. These external disturbances interact with the boundary layer, manifesting as fluctuations in the base flow that generate oscillatory modes whose amplitude, frequency, and phase define the initial conditions for instability waves that develop and eventually breakdown. This process, known as receptivity, converts external disturbances into instability waves whose initial amplitude plays a critical role in determining the onset and location of transition. In hypersonic boundary layers, key destabilising factors influencing transition in hypersonic flows include: \textbf{Freestream disturbances}: Acoustic waves, entropy fluctuations, and vorticity perturbations in the freestream can penetrate the boundary layer and trigger instability growth. \textbf{Surface roughness}: Imperfections on the surface can introduce localised perturbations that bypass classical instability mechanisms, leading to premature transition. \textbf{Shock interactions}: Shock-induced perturbations can introduce unsteady pressure and velocity fluctuations, modifying the growth of boundary layer instabilities. \textbf{High-temperature effects}: Vibrational excitation, dissociation, and recombination reactions in high-enthalpy flows alter the stability properties of the boundary layer, often enhancing instability growth. Because environmental disturbances are typically complex and broadband, the receptivity problem remains largely unsolved for compressible boundary layers, despite numerous efforts to address it.

% Transition begins with the appearance and amplification of disturbances from environmental sources such as freestream turbulence, acoustic noise, surface roughness, or shock interactions. These interact with the boundary layer, generating oscillatory modes whose amplitude, frequency, and phase define the initial conditions for instability waves that grow and ultimately breakdown. This \textit{receptivity process} governs how external disturbances are converted into instability waves, strongly influencing when and where transition occurs \citep{Morkovin1994}. In hypersonic boundary layers, key destabilising factors include freestream perturbations (acoustic, entropy, and vorticity fluctuations), surface roughness that introduces localised bypass routes, shock interactions that create unsteady pressure and velocity fields, and high-temperature effects like vibrational excitation and dissociation, all of which can significantly enhance instability growth \citep{vandereynde2015,Bitter2015}. Because these environmental disturbances are inherently complex and broadband, the receptivity problem remains largely unresolved for compressible boundary layers, despite extensive research \citep{dCerminara2017} Ma Zhong.

% In the most general case, low-amplitude oscillatory modes are triggered within the boundary layer and amplify through linear modal growth, followed by parametric instabilities and mode interactions that transfer energy between modes and lead to breakdown to turbulence; this transition pathway is depicted in \fref{figure:transitionmarkovin} as \textit{path A}. This path has been widely studied since the pioneering work of Tollmien and Schlichting, which established the well-known two-dimensional Tollmien--Schlichting instability waves. For higher disturbance levels, the superposition of perturbations may first exhibit algebraic transient growth before feeding into modal amplification (\textit{path B}). As the disturbance level increases further, transient growth can lead directly to parametric instabilities and mode interactions or, if sufficiently strong, bypass these stages entirely and transition straight to breakdown (\textit{paths C} and \textit{D}). Only \textit{paths A, B,} and \textit{C} are relevant for external hypersonic boundary layers, where disturbances grow through receptivity, transient growth, and modal amplification. \textit{Path D} mainly applies to internal flows with elevated turbulence levels, while \textit{path E} represents extreme forcing such as severe surface roughness or intense freestream turbulence where linear growth mechanisms are completely bypassed.


% Most of the groundwork done for  

% Ok, hear me out, most of the groundwork done for transition prediction was on incompressible flows which is not completely true but we say that anyway. The first studies about boundary layer stability was performed my Lees and Lin, who were quantifying the inflection point phenomenon, then came the Tollmein-Schliting waves, confirmed by further experiments by Schuberger or whatever their names are. It was not until Mack studies the stability of supersonic boudnary layer that we have a rigiorous definition of boundary layers and what they actually are. Mack performed stability analysis on supersonic base flows and characterised instability modes into numbers (first, second and third mode, so on). A detailed description of these instabilities is discussed in chapter 4, but the general gist of it is, first modes are promiment in subsonic incompressible flows, often associated with lower frequencies and the second mode is an acoustic trapped wave that is a feature of flows past supersonic regime. 



% The primary focus of the current thesis is to understand the modal decomposition of \textit{Path A}, 

% The most general transition process driven by freestream disturbances, known as Path A, involves low-amplitude disturbances that undergo exponential growth of eigenmodes. This linear amplification leads to parametric instabilities and mode interactions, which eventually cause breakdown to turbulence.  

% The most general transition process driven by freestream disturbances involves low-amplitude perturbations that undergo exponential growth of eigenmodes. This linear amplification leads to parametric instabilities and mode interactions, which effectively transfer energy between modes and eventually cause breakdown to turbulence. This pathway is illustrated in \fref{figure:transitionmarkovin} as \textit{path a}.

% In the most general case, low-amplitude oscillatory modes are triggered within the boundary layer and amplify through linear modal growth, followed by parametric instabilities and mode interactions that transfer energy between modes and lead to breakdown to turbulence. This transition pathway is depicted in \fref{figure:transitionmarkovin} as \textit{path a}. For slightly higher disturbance levels, the superposition of perturbations may exhibit algebraic transient growth before transitioning to modal growth; this transient growth mechanism, \textit{path b} which ultimately feeds into the same eigenmode amplification process as \textit{path a}, but is preceded by an initial phase of non-modal growth. For even higher disturbance levels, the flow may undergo transient growth that either leads directly to parametric instabilities and mode interactions (path C) or, if the disturbance level is sufficiently high, bypasses these stages entirely and transitions straight to breakdown (path D). Only paths A, B, and C are relevant for studying external hypersonic boundary layers, where disturbances grow through receptivity, transient growth, and modal amplification. Path D mainly applies to internal flows with elevated turbulence levels, where modal growth is bypassed altogether. Path E represents extreme disturbance environments — such as severe surface roughness or very high freestream turbulence — where linear mechanisms are entirely bypassed, leading directly to breakdown.




% Only paths A, B and C are
% sensible for the study of external hypersonic flows. Path D is primarily for internal flows
% at an elevated turbulence level. Path E represents the case of very large amplitude forcing
% (roughness and high freestream turbulence) where the growth of linear disturbances is
% bypassed




% The transition from a laminar to a turbulent boundary layer is a critical phenomenon in aerodynamics, with significant implications for skin friction, heat transfer, and aerodynamic drag. This process is initiated by amplifying flow disturbances, which may originate from various environmental sources such as freestream turbulence, acoustic noise, surface roughness, or shock interactions. The specific pathway to turbulence depends on the boundary layer's stability characteristics and the initial perturbations' nature and amplitude. A schematic overview of the various transition routes relevant to a flat plate boundary layer is shown in \fref{figure:transitionmarkovin} reproduced from \citep{Fedorov2011}.

% The transition from a laminar to turbulent boundary layer is a critical phenomenon in aerodynamics, with significant implications for skin friction, heat transfer


% \begin{figure}[h!]
%     \centering
%     % \includegraphics[width=13cm]{resources/figures/introduction/schematictempranges.png}
%     \includegraphics[scale=0.8]{example-image-c}
%     \caption{Markovin transition path to turbulence.}
%     \label{figure:transitionmarkovin}
% \end{figure}

% \begin{figure}[h!]
%     \centering
%     % \includegraphics[width=13cm]{resources/figures/introduction/schematictempranges.png}
%     \includegraphics[scale=0.8]{example-image-c}
%     \caption{Markovin transition path to turbulence.}
%     \label{figure:transitionmarkovin}
% \end{figure}





% According to the framework established by \cite{Morkovin}

% In incompressible and low-speed flows, the transition is typically governed by the amplification of two-dimensional Tollmien–Schlichting (TS) waves, which correspond to the \textit{first-mode} instability in Mack’s classification \citep{Mack1990}. These disturbances undergo exponential growth, eventually leading to nonlinear flow features such as $\Lambda$-vortices, hairpin vortices, and turbulent spots, culminating in fully developed turbulence. In contrast, hypersonic boundary layers are subject to additional instability mechanisms arising from compressibility effects, high-temperature-induced chemistry, and shock interactions. Key destabilising factors influencing transition in hypersonic flows include: \textbf{Freestream disturbances}: Acoustic waves, entropy fluctuations, and vorticity perturbations in the freestream can penetrate the boundary layer and trigger instability growth. \textbf{Surface roughness}: Imperfections on the surface can introduce localised perturbations that bypass classical instability mechanisms, leading to premature transition. \textbf{Shock interactions}: Shock-induced perturbations can introduce unsteady pressure and velocity fluctuations, modifying the growth of boundary layer instabilities. \textbf{High-temperature effects}: Vibrational excitation, dissociation, and recombination reactions in high-enthalpy flows alter the stability properties of the boundary layer, often enhancing instability growth.  


% A defining feature of hypersonic flows is the presence of \textbf{Mack modes}, or \textit{second-mode} instabilities, which arise from the interaction of trapped acoustic waves between the wall and the sonic line. The first theoretical studies of these compressible eigenmodes were conducted by Mack, who extended the classical linear stability theory (LST) developed by Orr and Sommerfeld. The second-mode instability in hypersonic boundary layers was first investigated numerically by \citet{Malik1990a}, who demonstrated its rapid amplification. Recent studies, such of \cite{Bitter2015} and \citet{miromiro2019}, have shown that dissociation effects in high-enthalpy regimes tend to destabilise second-mode instabilities, increasing their amplification rates. These findings also highlight the strong sensitivity of flow stability to variations in base flow conditions and wall properties.

% Building on these insights, recent advances in boundary-layer stability analysis have focused on computing second-mode $N$-factor envelopes using linear stability theory (LST), typically assuming frozen base flows \cite[p.~12-7]{miromiroThesis}. However, in post-shock regions—particularly near stagnation points—the boundary layer often approaches thermochemical equilibrium conditions \citep{Passiatore2021}, which challenges the validity of the frozen-flow assumption. This highlights the need to investigate how thermochemical state and wall boundary conditions influence stability characteristics. A deeper understanding of these effects is crucial for improving transition prediction models and informing control strategies in high-speed vehicle design.

% \section{Stability studies}
% Mention past stability studies that have gone into understanding the basic hypersonic regime

% \section{Numerical investigations}
% Talk about both temporal and spatial simulations for high enthalpy boundary layers


\newpage

\section{Recent Progress in Hypersonics}\label{section:introduction/recentprogress}
Recent progress in hypersonic aerothermodynamics has been driven by the renewed demand for sustained high-speed flight — not just for reentry capsules but also for maneuverable hypersonic glide vehicles (HGVs) and advanced air-breathing propulsion concepts. Historically, early vehicle designs relied on robust, heavily oversized thermal protection systems (TPS) and conservative margins because reliable prediction tools for transition and thermal loads were lacking. Empirical design rules, developed from limited flight data and wind tunnel tests, often led to over-designed heat shields and higher structural weights to account for uncertainties.

Advances in computational capabilities, experimental facilities, and materials science have steadily shifted this paradigm. Modern vehicles are now being designed for longer atmospheric flight times, tighter trajectory envelopes, and higher performance demands, all of which place greater emphasis on accurate thermal prediction. Material science developments — such as ultra-high-temperature ceramics, novel carbon–carbon composites, and oxidation-resistant coatings — have pushed operational temperature limits well beyond what was possible during earlier programs.

However, even with improved materials, the maximum surface temperatures achievable still fall short of the adiabatic wall temperatures that can occur on leading edges and stagnation points in sustained hypersonic flight. As a result, the use of cold wall conditions remains an integral design strategy to protect underlying structures and systems. This means that reliable modelling tools are needed to quantify not only the overall thermal loads but also the critical locations where transition from laminar to turbulent flow might occur earlier than expected due to the interaction of surface temperature and boundary layer instabilities.

Recent progress has focused on integrating advanced material capabilities with predictive design methods. Designers now face the challenge of leveraging these high-temperature materials efficiently while minimising mass and ensuring structural safety. Improved modelling tools that capture the multi-physics nature of hypersonic flows — including shock–boundary layer interactions, local heat fluxes, and realistic material behaviour under extreme conditions — are essential for closing the gap between material performance limits and real-world operational environments.

At the same time, the shift towards reusable vehicles and more frequent hypersonic missions has raised the stakes for design accuracy. Conservative margins are no longer sufficient when even small discrepancies in predicted heating rates can lead to costly overdesign or, worse, catastrophic failure. Better validated models allow engineers to make informed trade-offs between material selection, active cooling strategies, and flight trajectory design, ultimately enabling the next generation of hypersonic systems to operate safely and efficiently within ever tighter constraints.


% The complexity of modeling hypersonic turbulent flows relevant to practical applications stems from several unique physical phenomena that are either absent or far less pronounced at lower speeds. These include strong compressibility effects; the presence of intense shock waves that remain close to the vehicle body and the resulting shock/turbulent boundary layer interactions (SBLIs); extremely high temperatures due to severe viscous dissipation within turbulent boundary layers and strong shock waves, leading to chemical and thermal non-equilibrium and the associated turbulence–chemistry interactions (TCIs); as well as surface reactions such as ablation and the interaction of ablative particles with the flow (see Fig. 1).

% Intense aerothermodynamic heating, unsteady pressure loads, and high skin-friction drag are key global manifestations of these complex physical phenomena and are crucial considerations for engineering design. Furthermore, understanding the flow field—including the locations of shock waves, expansion fans, turbulent boundary and free-shear layers, and separated regions such as shock-induced recirculation bubbles, particularly under off-design conditions—is essential for the safe and successful operation of hypersonic vehicles. For example, during its final hypersonic flight, the X-15 aircraft experienced severe heating and structural damage due to shock waves from the engine cowl impinging on the pylon holding a dummy scramjet engine, as well as shock waves striking the vertical tail.

% Accurately predicting these complex hypersonic flow fields and the associated engineering quantities of interest (QoIs) remains a highly challenging task for the computational fluid dynamics (CFD) community.

% In recent years, significant progress has been made in hypersonic aerothermodynamics, with a particular emphasis on improving the prediction and physical understanding of complex flow phenomena, including boundary layer transition, shock–boundary layer interactions, and the effects of thermochemical nonequilibrium.

% Talk about: what physics happened, the vehicle and how they improved, how it changed and maybe melting points of TPS and why it is important from an experimental and modelling standpoint. Leading to experimental and modelling constraints section next. 


% Get some numbers about the TPS melting points. Get more numbers about

% \subsection{Instability growth}
% Considering a generic leading edge structure shown in \fref{figure:hypersonicwedge}, the transition derives from the flow conditions and any receptivity in the leading edge region, therefore, understanding the laminar boundary layer behaviour adjacent to the stagnation-point is important to understand the potential growth to turbulence transition. Understanding of the stagnation point flow in complete thermochemical non-equilibrium is
% important since this region states the conditions for the start of transition.


% \begin{figure}[h!]
%     \centering
%     \includegraphics[scale=0.36]{resources/figures/introduction/schematicshperecone.png}
%     \caption{Schematic of a sphere-cone geometry in a hypersonic flow field showcasing various regions in frozen, equilibrium and non-equilibrium conditions adapted from \citep{williams2021locally}.}
%     \label{figure:hypersonicwedge}
% \end{figure}

% The stagnation-point flow around a blunt body has been a subject of interest for since Hiemenz's numerical studies in 1911. The widely used approach for a dissociating gas in compressible flow was developed by \cite{FayRiddell1958}. The method provides an exactly self-similar solution to the Naiver-Stokes equations in dissociating gas to understand the stagnation point behaviour in the vicinity of a blunt geometry under various wall catalysis conditions. However, it should be mentioned that several authors (e.g. \citep{Lee2022} and \citep{Rocher2022}) have shown that these equations contain several approximations in their crude estimation of several crucial properties. 

% From an experimental standpoint, several empirical equations have been developed over the years to understand the heat-flux at the stagnation-point, most notably, Detra-Hidalgo \citep{Detra1961}, Detra-Kemp-Riddell, Brandis-Johnson \citep{Brandis2014} and Verant-Sagnier \citep{Sourgen2015}. More recently, \cite{Rocher2022} have extended the Fay-Riddell framework, incorporating semi analytical methods to include thermochemical equilibrium, while determining the surface species concentration in the stagnation point. However, the empirical nature of these correlations maybe inadequate for conditions outside the tested experimental regimes, prompting the requirement for more focused experiments to accurately capture the physics of the flow \citep{donaldson2021a}.  
% Concluding from the studies above, a sufficient self-similar framework incorporating chemical and vibrational non-equilibrium in various flow regions has been a vastly unexplored topic, providing scope for future research. 

% (Martin \& Candler, 1998, 1999, 2001) + Dyab \& Martin (2009)
% dCerminara2017,Cerminara2019,cerminarasandham2020



% \subsection{Experiments and Modelling}


Significant advances have been made in the last years in the study of high speed flows. Most detailed studies rely on CFD analyses, since carrying out physical experiments in the working conditions of interest is generally a costly or even infeasible task \citep{Bertin2006}. To get insights into the hypersonic regime, hyperveliocioty windtunnels have been designed to generate high speed flows in their working section. Unfortunately, the pwoer required to obtain the prescribed flow conditions is directly proporrtilnal to the flwo density, velocity, stagnationn enthalpy and cross-sectiopnal aread. Reproducing such severe conditions may cause overheating and result in structural damage; to avoid such problems, the test duration is often short. Moreover, instrumentation working efficiently at such extreme temperatures is also difficult to obtain. For all these reasons, hte cost of carrying the experiments are relatively high. However, there exist high ethalpy ground based facilities and in the flows we mention some of them. The NASA Hypersonic Tunnel Facility (HTF) originally designed to test nuclear thermal rockets, is capable of simulating up to Mach 7 at high enthalpy conditions. To simulate the effects of hypersonic conditions on TPS of re entry vehicles, the SCIROCCO and GHIBLI facilities of Plasma Wind Tunnel Complex at CIRA are able to reach Mach 12 and extremely high heat fluxes; the test chamber of SCIROCCO, in which the experimental measurements of pressure and temperature are carried out, is shown. In France, the PHEDRA wind tunnel simulates plasma gases of nitrogen air and methane and the ONERA F4 wint tuneell provides air flows at high enthlapy and pressure. The Plasmatron at Von Karman Institute is also a high enthalpy facility, which reaches temperatures up to 10000K; hot gas from the test chamber exits through a 700 kW heat exchanger and a 1050 kW cooling  system used. To overcome the aforementioned problematic, a number of hypersonic wind tunnels work at low enthalpy cryogenic conditions. The Purdue Mach 6 quiet tunnel, built by the end of 90s, is designed to achieve laminar nozzle-wall boundary layers to study natural transition to turbulence in cryogenic conditions. The Sandia Hypersonic wind tunnel accelerates the flow up to Mach 5 for an air mixture, and Mach 14 for a nitrogen mixture and the stagnation temperature is of the order of  1400K; this facility reproduces reentry vehicles or rockets flight conditions. A Mach number equal to 14 is reached by the AEDC Hypervelocity Wind Tunnel 9.


The standard experimental ground-based hypersonic facilities rarely operate at the conditions similar to that of a low altitude high-speed vehicle, as the facility power requirements to replicate the exact combination of velocity, flow density and therefore the stagnation enthalpy can be overwhelming to the current electrical grids \citep{passiatore2022}. Laminar and turbulent flows have received significant attention in past albeit at sub-orbital enthalpies, with approximations about recombination/dissociation reaction mechanisms. Several of these studies focusing on simplified turbulent hypersonic boundary layers with minimal chemical reactions were developed in the past in \cite{martin1998} , although it was not until \cite{duanmartin2011} that a rigorous differentiation between low-enthalpy and high-enthalpy calculations was provided, for a gas with little chemical activity neglecting thermal non-equilibrium. The studies have mostly considered temporal evolution of the boundary layer with little or no chemical activity. Therefore, little attention was directed towards understanding the nature at the extreme conditions of hypersonic flight \citep{Candler2019}.


\gls{DNS} of flow close to such extreme hypersonic regimes with thermochemical framework has recently appeared in the literature, with the first studies performed by \cite{marxen2013} on laminar boundary layers with instabilities suitable for hypersonic cruise conditions, whilst an extension study \citep{marxeniaccarinomagin2014} focused on non linear instability in the presence of finite rate chemistry without considering thermal non-equilibrium, concluding the dominant influence of chemical properties on primary disturbance amplitude, neglecting secondary instabilities during transition. Simulations of chemical non-equilibrium have received far greater attention as opposed to thermal non-equilibrium, with the \texttt{Mutation++} library \cite{Scoggins2020}, providing an extensive state of the art computation ability of physiochemical properties and has been widely used in the chemical studies concerning hypersonic flows. \cite{DiRenzo2020} introduced their \texttt{HTR} finite-difference solver for hypersonic aerothermodynamics for exascale computing, \cite{direnzourzay2021} used the routine to produce spatial evolution of such boundary layers to fully turbulent state. \cite{Passiatore2021} studied the complete spatially evolving turbulence transition process in thermochemical non-equilibrium for a quasi-adiabatic non-catalytic wall, where a second-mode transition was triggered with a strong amplitude and observed chemical activity to have little effect on the sub-harmonic transition. These studies were further extended to understand transition process in flat-plate boundary layer configurations with non-zero pressure gradients, i.e. oblique shock interactions by \citep{Volpiani2021} using another finite-difference code \texttt{CHARLIE}, and \cite{passiatore2023}, which is essentially an thermochemical extension of the shock wave interaction work of \cite{sandham2014} in hypersonic cold flow. The studies have concluded that mean velocities and turbulent intensity profiles are similar to those of cold hypersonic flows, in fact, the presence of varying species mass fractions caused a reduction in the turbulent fluctuations. At this point,  a complete understanding of chemical and thermal influences on modal stability for high Reynolds number, high Mach number flows is available in the literature. 



To date, various CFD solvers have been developed by some universities and scientific
research institutions for nonequilibrium flows in the near-continuum regime. These
include some proprietary solvers, such as NASA Langley code LAURA [58], NASA
Ames code DPLR

NASA Langley code \texttt{LAURA} \citep{Cheatwood1996}
NASA AMES code \texttt{DLPR} \citep{Candler1994}
\texttt{LeMANS} solver Michigan \citep{Scalabrin2005}
\texttt{US3D} Minnesota \citep{Nompelis2004}
\texttt{DLR-TAU} \citep{Mack2002}
\texttt{SU2-NEMO}, 
\texttt{HTR solver}


\section{Motivation}

% The primary motivation for this work is to advance the understanding of the different paths to transition in boundary-layer flows and to bridge the gap between the well-understood linear stages and the more complex nonlinear breakdown to turbulence. Accurate prediction of transition is critical for reliable estimates of aerodynamic drag, heat transfer, and vehicle performance, yet experimental investigations are often constrained by the energy demands and practical limitations of recreating realistic conditions, especially for high-speed flows. \gls{DNS} provides a powerful alternative by delivering high-fidelity, fully resolved data that can capture the entire transition process under controlled conditions. By using high-fidelity numerical studies, this work aims not only to clarify how freestream disturbances and receptivity mechanisms feed into the transition process, but also to generate benchmark data that can guide improved models and future experimental efforts.

% Guess bring everything together and summarise, why look into high enthalpy effects? why look at second mode? what is the need for a DNS solver, what are the issues with cold wall flows. 

Accurately predicting the transition from laminar to turbulent flow in hypersonic boundary layers remains one of the most challenging and consequential problems in high-speed aerothermodynamics. While linear stability theory (LST) has greatly improved our understanding of the initial growth of instabilities — particularly the second-mode (Mack mode) acoustic waves that dominate at higher Mach numbers — the detailed pathways through which these disturbances interact and break down into turbulence are still not fully understood.

Modal growth comparisons
This uncertainty is amplified in conditions relevant to modern hypersonic vehicles, where high temperatures and complex thermodynamic effects significantly alter the evolution of boundary layer instabilities. The consequences of transition prediction errors are substantial: premature transition can lead to localised heat loads well beyond what thermal protection systems are designed to handle, jeopardising structural integrity and mission safety. Advances in material science have pushed the operational temperature limits higher than ever before, but the need for surface cooling remains, which can further destabilise certain modes and complicate transition prediction.

Breakdown to turbulence things 

Traditional experimental facilities, such as hypersonic wind tunnels, have provided invaluable insights into transition mechanisms but face inherent limitations in replicating the full range of flight-relevant conditions — especially when it comes to capturing the coupled effects of high enthalpy and gas behaviour. Recent improvements in high-fidelity computational approaches, especially direct numerical simulation (DNS), have created new opportunities to investigate the multi-scale, nonlinear interactions that ultimately lead to turbulence. 

% Unlike empirical methods or lower-fidelity models, DNS can resolve the entire spectrum of instability growth and breakdown, making it an indispensable tool for uncovering the detailed mechanisms that govern transition in complex regimes.
writing a DNS solver to answer the questions

Despite this progress, key questions remain about how second-mode instabilities evolve in different thermodynamic environments and how they interact with other disturbance modes to form turbulent spots. A purely linear perspective cannot capture these interactions; therefore, a combined approach that integrates LST with DNS offers a more complete picture. Using LST to guide the identification of dominant modes and interpret their initial growth provides critical context for understanding how these instabilities develop into fully turbulent flow when all relevant physics are included.

How everything sums up together? 
This project aims to address these knowledge gaps by applying state-of-the-art DNS, complemented by LST analysis, to investigate the breakdown mechanisms of Mack modes in high-speed boundary layers under a range of thermal conditions. By clarifying how these instabilities grow, interact, and trigger transition, this work will contribute to the development of more reliable transition prediction methods. Such advances are essential for designing robust, efficient thermal protection systems and reducing uncertainty in the performance of next-generation hypersonic vehicles.


% \newpage

\subsection{Aims and objectives}\label{subsection:aimsandobjectives}

Following the motivation outlined above, this project aims to develop a rigorous mathematical and computational framework for analysing high-enthalpy hypersonic boundary layers under thermochemical non-equilibrium conditions, alongside a physical understanding of transition to turbulence. Ionisation effects are not considered in this study. The research first seeks to extend self-similar boundary-layer theory, which has received relatively little attention in the low-altitude, high-enthalpy regime where multiple species and nonequilibrium effects play a significant role. While most \gls{LST} studies in the literature assume a frozen freestream base flow~\citep{miromiroThesis}, the post-shock region near the stagnation point is typically closer to local thermochemical equilibrium, which remains poorly understood. Furthermore, the project investigates the transition mechanisms in high-Mach-number, high-Reynolds-number conditions under varying wall boundary conditions.

To address these challenges, part of the work focuses on extending the in-house finite difference compressible solver OpenSBLI to incorporate detailed thermal and chemical nonequilibrium models with appropriate shock-capturing techniques. The project also aims to clarify how extreme nonlinear interactions drive transition to turbulence under these conditions. Taking into account the current state of the art, the following objectives are proposed to address this critical gap in our understanding of transition on a wedge and the leading edge of a blunt body in thermochemical nonequilibrium.




% Furthermore, the project intends to understand the bypass transition in the realms of hypersonic flow under different wall conditions. In particular, the aims to understand the nature of the flow in transition region to turbulence in presence of extreme non-linear effects. 
% Part of the work is dedicated to extend the current in-house compressible solver to include thermal and chemical nonequilibrium with appropriate models and shock chapturing techniques. 


% \begin{enumerate}[topsep=0pt,itemsep=-1ex,partopsep=1.5ex,parsep=3.5ex]
%     \item To extend the analytical set of Fay-Riddell equations to include multiple species,  thermochemical non-equilibrium, and Park's two-temperature model.
%     % \item To perform DNS simulations of a flow around a cylinder by extending OpenSBLI for dissociating gas incorporating equations \sref{section:literature/navier-stokes} for necessary validations. 
%     \item Extend the in-house compressible solver OpenSBLI to include thermochemical framework with chemical and vibrational non-equilibrium, including validation for flow around a cylinder to understand the stagnation point behaviour. 
    
%     % \item To perform DNS simulations of flow around a cylinder to understand the stagnation point behaviour under equilibrium freestream, including validation for .
%     % \item To accurate incorporate DNS methods with suitable shock-capturing scheme to simulate the laminar-turbulent transitions on a flat plate. 
    
%     % \item To accurate incorporate DNS methods with suitable shock-capturing scheme to simulate the laminar-turbulent transition. 
%     % \item 
%     \item To investigate the effects of LST and its application in e\textsuperscript{N} method to accurately predict the behaviour of instability modes and their sensitivities to the high-enthalpy base flow inaccuracies.

    
%     \item Understand the mechanics of temporally and spatially developing turbulent flow region, using DNS of spreading turbulent spots and wedge of developing turbulent flow. 

%     % \item 
%     % \item Apply e\textsuperscript{N} method to the modal analysis to understand the sensitives of turbulence transition.
%     % \item 
% \end{enumerate}

% The physics of hypersonic flows differs vastly from the subsonic and supersonic counterparts. Most vehicles in hypersonic flows encounter strong shock-wave compression resulting in high-enthalpy gas environment that almost always associates with non-equilibrium thermodynamic and  chemical phenomena. These high-enthalpy conditions often render the calorically perfect gas models moot, consequently a multi-species model to capture the chemically reacting atmosphere is required. Further complexities such as shock-waves, in addition of these chemical reactions, add more challenges to accurately predict the loads on the vehicle.  

% still a foreigner concept along with the understanding of the mechanisms to transition, despite the great deal of attention received by them over the years. 

% The outline of the current report is as follows: necessary theoretical framework for the project is suggested in Chapter 2, whilst the numerical framework along with the simulation method is ascribed in Chapter 3. Chapter 4 covers the self-similar studies for both stagnation region and numerical inflow conditions with thermochemistry. Currently extensions to the compressible Navier-Stokes DNS solver to include thermochemistry are mentioned in Chapter 5 along with studies on cylinder flow, flat plate simulations and preliminary transition results. Finally, conclusions of the current work are drawn along with necessary framework for the thesis completion. 

% \newpage

\section{Thesis outline}
The structure of this report is as follows: \Cref{chapter:governing_equations} presents the necessary theoretical framework and governing equations used to describe the fundamental flow physics, including the thermochemical processes relevant to high-enthalpy boundary layers. \Cref{chapter:numerical_methods} introduces the \gls{DNS} solver and details the numerical methods employed, along with representative validation cases that demonstrate the solver’s ability to resolve the complex features of transitional hypersonic flows. \Cref{chapter:disturbance_growth} examines the propagation and amplification of small disturbances leading to second-mode breakdown in thermally excited boundary layers, drawing systematic comparisons between \gls{LST} and \gls{DNS} results to highlight the role of vibrational nonequilibrium and real gas effects. Building on this, \Cref{chapter:breakdown} focuses on isolating and analysing Mack modes to clarify the mechanisms by which these dominant instability modes undergo complete breakdown and trigger transition under various thermally excited conditions. Finally, \Cref{Chapter: Conclusions} summarises the main findings, discusses their broader implications for hypersonic transition prediction and \gls{TPS} design, and provides recommendations for future research directions. \gls{CPG}, \gls{TEQ},\gls{TNE}

% The outline of the report is as follows: necessary theoretical framework for the project is suggested in \Cref{chapter:numerical_methods}, and the numerical framework employed for the theoretical framework along with the simulation method is ascribed in Chapter 3; which also includes the implementation of the catalytic wall boundary condition on a flow around a cylinder. Chapter 4 covers the self-similar studies for stagnation point flow past a dissociated shock layer, along with suggestions for inflow boundary conditions for compressible thermochemical simulations. Current endeavours to extend the compressible Navier-Stokes DNS solver OpenSBLI to include thermochemistry is mentioned in Chapter 5, which provides preliminary studies on flat plate laminar and transitional boundary layers in both cold and high-enthalpy supersonic flows. Finally, conclusions of the current work are drawn in Chapter 6 along with necessary future framework for the thesis completion. 